\documentclass[10pt,a4paper]{article}
\usepackage[utf8]{inputenc}
\usepackage[T1]{fontenc}
\usepackage{lmodern}
\IfFileExists{french.ldf}{\usepackage[french]{babel}}{}
\usepackage[margin=19mm,headheight=14pt]{geometry}
\usepackage{amsmath,amssymb,graphicx,booktabs,array,xcolor}
\usepackage{microtype,fancyhdr,tcolorbox}
\AtBeginEnvironment{itemize}{\setlength{\itemsep}{3pt}\setlength{\parskip}{1pt}}
\tcbuselibrary{breakable}
\usepackage[colorlinks=true,linkcolor=blue!45!black,urlcolor=blue!60!black,citecolor=blue!45!black]{hyperref}
\renewcommand{\contentsname}{Sommaire}
\renewcommand{\refname}{Références}
\definecolor{navy}{HTML}{183047}
\definecolor{teal}{HTML}{167D9A}
\definecolor{gold}{HTML}{DBA52D}
\definecolor{pale}{HTML}{F2F6F9}

\pagestyle{fancy}
\fancyhf{}
\fancyhead[L]{\small\textcolor{navy}{Cu\textsubscript{3}Au : ordre et désordre chimique}}
\fancyhead[R]{\small Projet Monte-Carlo}
\fancyfoot[L]{\small Réseau CFC rigide / ensemble canonique}
\fancyfoot[R]{\thepage}
\setlength{\parindent}{0pt}
\setlength{\parskip}{3pt}
\setlength{\emergencystretch}{2em}
\setlength{\textfloatsep}{9pt plus 2pt minus 2pt}
\setlength{\floatsep}{7pt plus 2pt minus 2pt}
\setlength{\intextsep}{8pt plus 2pt minus 2pt}
\setlength{\abovecaptionskip}{4pt}
\setlength{\belowcaptionskip}{2pt}
\setcounter{topnumber}{3}
\setcounter{bottomnumber}{2}
\setcounter{totalnumber}{4}
\renewcommand{\topfraction}{.90}
\renewcommand{\bottomfraction}{.85}
\renewcommand{\textfraction}{.08}
\renewcommand{\floatpagefraction}{.75}
\setcounter{tocdepth}{2}
\newtcolorbox{oral}[1]{breakable,colback=pale,colframe=teal,boxrule=.6pt,
  title={\textbf{Approfondissement pédagogique : #1}},fonttitle=\footnotesize,coltitle=white,
  fontupper=\small,top=2pt,bottom=2pt,before skip=5pt,after skip=5pt}
\newtcolorbox{attention}{breakable,colback=gold!8,colframe=gold!80!black,boxrule=.6pt,
  title={Point d'interprétation},fonttitle=\small}
\newtcolorbox{algocode}[1]{breakable,colback=pale,colframe=navy,boxrule=.6pt,
  title={#1},fonttitle=\bfseries,coltitle=white,
  before={\par\smallskip},after={\par\smallskip}}
\newcommand{\Au}{\mathrm{Au}}
\newcommand{\Cu}{\mathrm{Cu}}
\newcommand{\kb}{k_{\mathrm B}}
\newcommand{\fig}[3]{\begin{figure}[htbp]\centering
\includegraphics[width=#2\linewidth,height=.42\textheight,keepaspectratio]{figures/#1.pdf}
\begingroup\small\caption{#3}\label{fig:#1}\endgroup\end{figure}}
\input{auto_results.tex}
\begin{document}
\hypersetup{pageanchor=false}
\begin{titlepage}
\vspace*{5mm}
{\color{teal}\rule{\linewidth}{2pt}}
\vspace{9mm}
{\Large\textcolor{teal}{Modélisation et simulation numérique}\par}
\vspace{8mm}
{\Huge\bfseries\color{navy} Transition ordre--désordre de Cu\textsubscript{3}Au\par}
\vspace{5mm}
{\large Étude Monte-Carlo sur réseau rigide : méthode, résultats et limites\par}
\vspace{8mm}
\begin{tcolorbox}[colback=pale,colframe=navy,boxrule=.8pt]
\textbf{Système étudié :} $\Lbox^3$ mailles, $N=\Natoms$ sites,
$c_{\Au}=1/4$, interactions de premiers et seconds voisins.\\[2mm]
\textbf{Résultat exact :} $u_0=-3.813\ \mathrm{eV/atome}$.\\[2mm]
\textbf{Maximum de la production canonique L6 :} $T^*_{C_V}=\Tpeak\ \mathrm K$ sur la grille.\\[2mm]
\textbf{Diagnostics de cette production près du maximum :} \DiagnosticStatus.
\IfFileExists{bonus_sizes.tex}{\par\small
\textbf{Complément achevé :} comparaison croisée $L=6,9,12$ en
section~\ref{sec:tailles}. Les diagnostics L9 et L12 près du pic restent
insuffisants ; le maximum canonique L12 observé vaut 1052.5 K.
Les cycles chauffage--refroidissement sont achevés pour $L=\CompletedCycleSizes$.}{}
\end{tcolorbox}
\vfill
\textbf{Traçabilité.} Tous les résultats numériques, tableaux et figures de ce
rapport proviennent des fichiers du dossier de résultats associé. Le code
recalcule les résultats à partir des paramètres du modèle ; aucune valeur
de transition n'est imposée à la simulation.
\vspace{8mm}
\end{titlepage}
\hypersetup{pageanchor=true}
\tableofcontents
\newpage

\section*{Cadre du projet et conventions}
\addcontentsline{toc}{section}{Cadre du projet et conventions}
Le projet étudie le désordre \emph{chimique} d'un alliage Cu$_3$Au. Les sites
du réseau ne se déplacent pas : l'état d'un site est seulement son occupation
par Au ou Cu. La composition et le volume sont fixés. Les températures sont
en kelvins, les énergies en eV et $\kb=8.617333262\times10^{-5}\ \mathrm{eV/K}$.
Dans le code, $s_i=1$ désigne Au et $s_i=0$ Cu. Une paire non orientée est
comptée une seule fois.
Le paramètre $a$ est le côté de la maille conventionnelle. Les longueurs
sont données en unités de $a$ ; une conversion en nanomètres demanderait
une valeur physique de $a$, qui n'est pas imposée dans cette simulation.

\begin{center}
\begin{tabular}{lrrr}
\toprule
Coquille & Distance & Coordination & Énergies de paire (eV)\\
\midrule
Premiers voisins & $a/\sqrt2$ & $12$ & $J_{\Au\Au}=-0.635$\\
 & & & $J_{\Cu\Cu}=-0.582$, $J_{\Au\Cu}=-0.689$\\
Seconds voisins & $a$ & $6$ & $I_{\Au\Au}=I_{\Cu\Cu}=0$\\
 & & & $I_{\Au\Cu}=+0.0063$\\
\bottomrule
\end{tabular}
\end{center}

\subsection*{Repères de vocabulaire}
Un \textbf{site} est une position fixe du réseau ; son occupation est Au ou Cu.
Une \textbf{maille} est le motif géométrique répété dans l'espace. Les quatre
\textbf{sous-réseaux}, repérés par $\mu=0,1,2,3$, distinguent les positions
de base de la maille CFC. Une \textbf{coquille} regroupe les voisins à une
même distance d'un site choisi. Une \textbf{configuration} donne l'espèce
sur chacun des sites à un instant du calcul.

Une \textbf{réplique} est une configuration qui peut voyager entre les
températures. Une \textbf{chaîne} désigne ici une campagne indépendante
avec sa graine et son ensemble de répliques (\Ntemps\ dans la production L6).
Dans le cycle séquentiel, une chaîne est au contraire une seule trajectoire
transmise d'une température à la suivante. Un \textbf{balayage} effectue
$N$ propositions Au--Cu par température, refus compris. L'\textbf{équilibration}
est la phase initiale écartée des moyennes ; la \textbf{production} est la
phase dont on utilise les mesures pour estimer les grandeurs thermiques.
Un \textbf{domaine} est une région spatiale partageant une variante d'ordre ;
ce n'est ni une chaîne ni un sous-réseau. Le symbole $\mu$ est ici un
indice géométrique de sous-réseau, pas un potentiel chimique : aucun terme
de potentiel chimique n'entre dans le test de Metropolis à composition fixée.
Enfin, $L1_2$ est le nom de la structure ordonnée ; $L=12$ est la taille
de la boîte. Les abréviations L6, L9 et L12 du rapport désignent
respectivement $L=6,9,12$, et non trois structures cristallines différentes.

\subsection*{Méthodologie : quatre expériences numériques distinctes}
\addcontentsline{toc}{subsection}{Méthodologie et portée des mesures}
Le Hamiltonien, la composition et les conditions périodiques sont communs,
mais les objectifs et les données ne sont pas interchangeables.
Les résultats de référence des questions 1--9 proviennent de la production
canonique principale L6. La comparaison achevée avec L9 et L12 est regroupée en
section~\ref{sec:tailles}.
Les expériences du bonus sont identifiées séparément ; leurs nombres de
balayages et leurs diagnostics ne remplacent pas ceux de la campagne principale.
\begin{center}\small
\begin{tabular}{>{\raggedright\arraybackslash}p{.23\linewidth}>{\raggedright\arraybackslash}p{.31\linewidth}>{\raggedright\arraybackslash}p{.35\linewidth}}
\toprule
Campagne $L=6$ & Protocole par chaîne et température & Données et usage légitime\\\midrule
Canonique principale & 46 températures, 2 chaînes ; 6000 balayages écartés,
30000 mesurés, pas 10 ; échanges de répliques tous les 5 balayages &
Énergie, ordre, fluctuations et $T^*(L)$ ; états atomiques finaux seulement.\\
Cycle thermique & 46 températures dans chaque sens, 2 trajectoires ;
6000 + 30000 balayages, pas 10 ; aucun échange de répliques &
Mémoire du parcours à budget fini ; un état final par palier pour le GIF.\\
Trempe isotherme & Une trajectoire, consigne $4000\to1000$ K ;
6000 balayages, états tous les 20, aucun transitoire retiré &
Relaxation hors équilibre ; pas de distribution de temps de nucléation.\\
Corrélations canoniques & 14 températures, 2 chaînes ; 3000 + 12000 balayages,
états tous les 20 ; échanges de répliques tous les 5 &
600 états par chaîne et température ; spectres et longueurs conditionnels aux diagnostics.\\
\bottomrule
\end{tabular}\end{center}
Les productions canoniques supplémentaires utilisent le même modèle :
$L=9$ conserve les 46 températures et le budget L6 ; $L=12$ utilise
90 températures, 12000 balayages écartés et 60000 de production par chaîne.
Les captures spatiales de L9 et L12 sont espacées de 100 et 200 balayages.
Ces productions simultanées en température ne sont pas des cycles :
les tailles dont les branches séquentielles sont entièrement achevées
dans cette version sont $L=\CompletedCycleSizes$.

Une graine indépendante crée une nouvelle trajectoire ; les répliques
d'une même chaîne restent couplées par leurs échanges. Les mesures sont
classées par température, après les éventuels échanges. Écarter un nombre
fixé de balayages ne démontre pas l'équilibre : les tests de correction,
la stationnarité observée et l'exploration des variantes répondent à des
questions différentes. Une série gelée peut réussir un test de stabilité
tout en manquant les excitations rares. Les variances thermiques sont
calculées dans chaque chaîne ; une divergence entre chaînes ne doit pas
être convertie en chaleur spécifique.

\textbf{Trois niveaux de certitude.} La géométrie, les différences d'énergie
et les références à $T=0$ sont exactes pour le Hamiltonien donné.
Les moyennes et intervalles Monte-Carlo sont des estimations conditionnelles
à la durée, aux états visités et à la grille. La comparaison à un alliage
réel ajoute une erreur de modèle, qui n'est pas contenue dans les barres
statistiques. Le calcul en micro-eV évite la dérive d'accumulation ; il ne
rend pas les paramètres physiques plus précis que ceux de l'énoncé.

\subsection*{Comment lire les résultats sans mélanger les protocoles}
\begin{itemize}
\item \textbf{Questions 1--9 : référence L6.} Sauf mention contraire,
les figures, les nombres et les diagnostics de ces questions concernent
la production canonique de 864 sites. Les résultats des grandes tailles
ne remplacent pas silencieusement cette référence.
\item \textbf{Bonus : quatre questions complémentaires.} La partie 1
explique le cycle et les GIF L6 ; la partie 2 définit les corrélations
spatiales ; la partie 3 compare les productions canoniques L6/L9/L12 ;
la partie 4 compare leurs cycles séquentiels achevés. Les GIF restent L6.
\item \textbf{Toujours lire valeur, protocole et réserve ensemble.}
$q$ mesure une amplitude globale, $\alpha_1,\alpha_2$ deux voisinages,
$\xi/a$ une portée spectrale. Le temps d'autocorrélation $\tau$ mesure
une dépendance entre mesures successives, en mesures ou en balayages :
ce n'est pas une longueur ni un temps physique en secondes.
\end{itemize}

\section{Question 1 -- Réseau, énergie et algorithme de Metropolis}
\subsection{Construction exacte de la géométrie}
Une maille conventionnelle porte les quatre positions
\[
\mathbf b_0=(0,0,0),\quad
\mathbf b_1=(0,\tfrac12,\tfrac12),\quad
\mathbf b_2=(\tfrac12,0,\tfrac12),\quad
\mathbf b_3=(\tfrac12,\tfrac12,0).
\]
Les sites sont $\mathbf r/a=(i,j,k)+\mathbf b_\mu$, avec
$i,j,k=0,\ldots,L-1$. Ainsi $N=4L^3$, $N_{\Au}=L^3$ et $N_{\Cu}=3L^3$.
La boîte demandée correspond à $L=6$, donc $864$ sites dont $216$ Au.

Le changement de coordonnées $(X,Y,Z)=2\mathbf r/a$ permet de travailler
avec des entiers. Les sites CFC vérifient
\[
X,Y,Z\in\{0,\ldots,2L-1\},\qquad X+Y+Z\text{ pair}.
\]
En effet, les parités possibles sont $(0,0,0)$, $(0,1,1)$, $(1,0,1)$ et
$(1,1,0)$, soit exactement les quatre positions de base. Pour $L=6$, les
coordonnées vont de $0$ à $11$. La coordonnée $12$ est identifiée à $0$.

Les déplacements de premiers voisins sont
\[
(\pm1,\pm1,0),\quad(\pm1,0,\pm1),\quad(0,\pm1,\pm1),
\]
et ceux de seconds voisins sont
\[
(\pm2,0,0),\quad(0,\pm2,0),\quad(0,0,\pm2).
\]
Toutes les coordonnées sont réduites modulo $2L$. Le tableau des voisins
est calculé une fois et reste fixe. Les tailles $L<3$ sont refusées dans cette
implémentation pour éviter les ambiguïtés d'images périodiques multiples.

\fig{01_reseau}{.95}{La maille $L1_2$ et le partage des sites de bord. Les
positions sont celles du réseau CFC, quelle que soit l'occupation chimique.}

\subsection{Énergie et tendance à l'ordre}
\fig{13_coquilles}{.95}{Les deux coquilles autour d'un Au dans la référence
ordonnée $L1_2$. Les douze premiers voisins sont Cu et les six seconds
voisins sont Au. Les coquilles sont définies par leurs distances, même si
leurs occupations changent pendant la simulation.}

L'énergie est
\begin{equation}
E=\sum_{\{i,j\}\in\mathcal P_1}J_{\sigma_i\sigma_j}
+\sum_{\{i,j\}\in\mathcal P_2}I_{\sigma_i\sigma_j}.
\end{equation}
Si l'on utilise les listes de voisins orientées, il faut un facteur $1/2$.
On obtient $Nz_1/2=5184$ paires de premiers voisins et $Nz_2/2=2592$ paires
de seconds voisins pour $N=864$.

Pour comparer les préférences chimiques, on définit
\[
w_1=J_{\Au\Au}+J_{\Cu\Cu}-2J_{\Au\Cu}=0.161\ \mathrm{eV}.
\]
Deux contacts Au--Cu coûtent donc $0.161$ eV de moins qu'un contact Au--Au
et un contact Cu--Cu. Au second voisinage, chaque paire Au--Cu coûte
$0.0063$ eV, tandis qu'une paire identique coûte zéro. L'ordre $L1_2$
satisfait ces deux tendances.

\subsection{Échanges à composition fixée et variation locale}
Un mouvement échange un Au et un Cu, choisis uniformément parmi les paires
d'espèces différentes. Cela conserve exactement la stœchiométrie.
Pour des sites $i$ et $j$, initialement Au et Cu,
\begin{align}
\Delta E={}&\sum_{n=1}^2\sum_{k\in V_n(i)\setminus\{j\}}
[U_n(\Cu,\sigma_k)-U_n(\Au,\sigma_k)]\notag\\
&+\sum_{n=1}^2\sum_{k\in V_n(j)\setminus\{i\}}
[U_n(\Au,\sigma_k)-U_n(\Cu,\sigma_k)],
\end{align}
où $V_n(i)$ est la liste des voisins et $U_1=J$, $U_2=I$ les énergies de paire.
Si les sites sont voisins, leur liaison mutuelle
reste Au--Cu : sa variation est nulle. Elle ne doit pas être recalculée comme
si une seule extrémité changeait.

Le noyau rapide exploite aussi l'identité démontrée à la question 5.
Si $a_i,a_j$ sont les nombres de voisins Au de première coquille, et
$b_i,b_j$ ceux de seconde coquille, en excluant le partenaire échangé,
\begin{equation}
\boxed{\Delta E=0.161(a_j-a_i)+0.0126(b_i-b_j)\quad\text{eV}.}
\end{equation}
Les énergies sont stockées en micro-eV entiers. La somme locale est ainsi
exacte, sans accumulation d'erreurs d'arrondi. Les exponentielles
d'acceptation restent naturellement évaluées en nombres réels.

\subsection{Balance détaillée}
Avec $\beta=1/(\kb T)$, la mesure canonique sur les configurations de composition fixée est
\[
\pi_T(C)=Z^{-1}e^{-E(C)/(\kb T)}.
\]
Comme la proposition d'échange est symétrique, l'acceptation de Metropolis est
\begin{equation}
\boxed{A(C\to C')=\min[1,e^{-\Delta E/(\kb T)}].}
\end{equation}
On a $A(C\to C')/A(C'\to C)=e^{-\beta\Delta E}=\pi(C')/\pi(C)$,
ce qui établit la balance détaillée. Accepter certains mouvements défavorables
est nécessaire pour reproduire les excitations thermiques.

Un balayage représente ici $N$ propositions \emph{Au--Cu}. Les propositions
triviales entre deux espèces identiques sont retirées avant le comptage.
Cette convention modifie l'unité de temps Monte-Carlo, mais pas la mesure
d'équilibre. Les échanges à longue distance sont un outil d'échantillonnage ;
ils ne décrivent pas une cinétique réaliste de diffusion par lacunes.

\fig{14_organigramme}{.72}{Organisation du calcul : proposition, décision,
comptage des tentatives, fin de balayage, échanges puis mesures. Le dessin
résume les étapes ; les boucles complètes sont données en annexe.}

\begin{oral}{Question 1 -- Pourquoi cet algorithme et dans quel ordre ?}
\begin{itemize}
\item \textbf{Échanger pour conserver la composition.} Changer un seul Au
en Cu ferait sortir de l'ensemble de configurations demandé. En choisissant
un Au et un Cu, on garde exactement $N/4$ Au. La probabilité de proposer
l'échange inverse est identique : l'acceptation dépend seulement du rapport
des poids de Boltzmann. Le label géométrique $\mu$ reste celui du site ;
seule son occupation change.
\item \textbf{Calculer avant de modifier.} On évalue $\Delta E$ sur
l'état courant en examinant les liaisons touchées, puis on applique
Metropolis. Si $\Delta E>0$, accepter parfois évite une simple descente
vers un minimum. À 1000 K, un défaut de 0.5586 eV serait accepté avec une
probabilité $\exp[-0.5586/(k_B1000)]\simeq0.00153$ pour cette proposition.
Un refus compte comme une tentative et conserve le temps de séjour de l'état.
\item \textbf{Mesurer après une durée fixée.} Un balayage contient $N$
propositions, pas $N$ acceptations. Les échanges de répliques éventuels
viennent après les mouvements, puis la mesure est rangée à température
fixe. Mesurer uniquement les changements acceptés surreprésenterait les
états faciles à quitter. Les voisins sont précalculés car la géométrie
ne change jamais ; le facteur $1/2$ intervient dans les sommes
où une même paire est vue depuis ses deux extrémités.
\end{itemize}
\end{oral}

\section{Question 2 -- Ordre à courte portée : paramètres de Cowley}
Pour une référence Au, l'énoncé définit
\[
\alpha_n=1-\frac{p_n^{\Cu\mid\Au}}{c_{\Cu}},\qquad c_{\Cu}=\frac34.
\]
Si $B_n$ est le nombre de paires Au--Cu non orientées de la coquille $n$,
chaque paire possède exactement une extrémité Au. On obtient
\[
p_n^{\Cu\mid\Au}=\frac{B_n}{N_{\Au}z_n},\qquad
\boxed{\alpha_n=1-\frac{B_n}{N_{\Au}z_nc_{\Cu}}.}
\]
Pour la boîte de l'énoncé,
\[
\alpha_1=1-\frac{B_1}{1944},\qquad
\alpha_2=1-\frac{B_2}{972}.
\]
\begin{tcolorbox}[breakable,colback=pale,colframe=teal,title={Du comptage à la probabilité puis à Cowley}]
À première coquille, les $216$ Au donnent $216\times12=2592$ positions
voisines examinées. Si $1944$ de ces positions portent Cu, alors
\[
p_1^{\Cu\mid\Au}=\frac{1944}{2592}=0.75,
\qquad \alpha_1=1-\frac{0.75}{0.75}=0.
\]
La première division mesure une probabilité conditionnelle ; la seconde
la compare à la concentration globale $c_{\Cu}=0.75$.
Si les $2592$ positions portent Cu, $p_1^{\Cu\mid\Au}=1$ et
$\alpha_1=-1/3$. À seconde coquille, on examine $216\times6=1296$
positions : si aucune ne porte Cu, $p_2^{\Cu\mid\Au}=0$ et $\alpha_2=1$.
Le code effectue ce comptage autour de tous les Au, à chaque mesure,
puis moyenne les paramètres obtenus sur la production.
\end{tcolorbox}

Dans $L1_2$, les premiers voisins d'un Au sont tous Cu, donc
$\alpha_1=1-1/(3/4)=-1/3$. Les seconds voisins sont tous Au, donc
$\alpha_2=1$.

Un paramètre négatif indique un excès de voisins d'espèce opposée. Un
paramètre positif indique leur déficit. Zéro est la référence d'un mélange
indépendant pour cette coquille dans la limite d'un grand système.
La nullité de deux paramètres ne prouverait pas, à elle seule, l'absence
de toute corrélation à d'autres distances.

À composition fixée dans une boîte finie, conditionner la référence à être
Au laisse $N_{\Cu}$ Cu parmi $N-1$ autres sites. Ainsi
\[
\left\langle\alpha_n\right\rangle_{T\to\infty}
=1-\frac{N_{\Cu}/(N-1)}{N_{\Cu}/N}
=\boxed{-\frac1{N-1}}.
\]
Pour $N=864$, cette correction vaut environ $-0.00116$.

\fig{12_cowley}{1}{Évolution des paramètres de Cowley $\alpha_1$ et
$\alpha_2$ avec la température. Les barres représentent une erreur-type.
Les lignes horizontales indiquent l'ordre parfait $L1_2$.}

À basse température, les courbes rejoignent les valeurs de l'ordre parfait :
$\alpha_1=-1/3$ et $\alpha_2=1$. Lorsque la température augmente,
les deux paramètres s'éloignent de ces valeurs et se rapprochent de zéro.
Un $\alpha_1$ encore négatif au-dessus de la zone de perte d'ordre global
signifie que les voisins Cu restent favorisés autour d'Au : l'ordre local
peut donc persister après la disparition de l'ordre à longue portée.
La limite aléatoire correspond à $T\to\infty$ ; les mesures à
température finie ne sont pas contraintes à l'atteindre.

\begin{oral}{Question 2 -- Des voisins comptés à la probabilité puis à Cowley}
\begin{itemize}
\item \textbf{Identifier la population de référence.} On ne compte les
voisins qu'autour des sites actuellement Au. Pour la première coquille L6,
il existe $216\times12=2592$ places de voisinage. Si 2160 portent Cu,
$p_1^{\Cu|\Au}=2160/2592=5/6$. C'est une probabilité conditionnelle,
comprise entre 0 et 1 ; ce n'est pas encore $\alpha_1$.
\item \textbf{Comparer à la concentration.} Avec $c_{\Cu}=3/4$,
$\alpha_1=1-(5/6)/(3/4)=-1/9$. Le signe négatif indique que Cu est
plus fréquent autour d'Au que dans la composition globale. Dans l'ordre
parfait, les douze voisins sont Cu, donc $\alpha_1=-1/3$ ; dans la
seconde coquille, aucun n'est Cu, donc $\alpha_2=1$. Deux coquilles
peuvent ainsi avoir des signes opposés sans contradiction.
\item \textbf{Distinguer ordre local et ordre de la boîte.} Cowley ignore
quel sous-réseau porte globalement les Au ; il est invariant entre les
quatre variantes. Des voisinages préférentiels peuvent survivre lorsque
les domaines se compensent dans l'ordre global. Une moyenne proche de
zéro est une référence de mélange, avec correction canonique
$-1/(N-1)$ à température infinie ; à température finie, les interactions
empêchent en général d'atteindre exactement cette limite.
\end{itemize}
\end{oral}

\section{Question 3 -- Ordre à longue portée et variantes équivalentes}
Le sous-réseau des sommets comprend $N/4$ sites. En notant
$n_{\Au}^{(0)}$ son nombre d'Au,
\[
p_0=\frac{n_{\Au}^{(0)}}{N/4},\qquad
\boxed{\eta=\frac{p_0-1/4}{3/4}}.
\]
Pour $N=864$, cela s'écrit $\eta=(n_{\Au}^{(0)}-54)/162$.
Dans la variante de référence, $p_0=1$ et $\eta=1$. Dans un mélange
aléatoire, $\langle p_0\rangle=1/4$ et $\langle\eta\rangle=0$.

Il existe quatre variantes équivalentes : Au peut occuper chacun des quatre
sous-réseaux CFC. Dans une autre variante parfaitement ordonnée, $p_0=0$
et $\eta=-1/3$. Pour éviter une mauvaise interprétation, le code conserve
le paramètre du sujet et calcule aussi
\begin{equation}
q=\sqrt{\frac43\sum_{\mu=0}^3(p_\mu-1/4)^2},
\qquad \sum_{\mu=0}^3p_\mu=1.
\end{equation}
On a $q=1$ dans les quatre variantes et $q=0$ pour les occupations moyennes
uniformes. Mais $\langle q\rangle$ est légèrement positif dans une boîte finie
désordonnée, car il mesure une amplitude de fluctuations.
Inversement, un $q$ faible ne prouve pas une occupation indépendante :
plusieurs domaines de variantes différentes peuvent réduire l'ordre global
tout en conservant des voisinages localement ordonnés.

\fig{15_variantes}{1}{Les quatre variantes parfaitement ordonnées sont
équivalentes : seul le sous-réseau qui porte les Au change. $\eta$ dépend
du sous-réseau de référence, tandis que $q=1$ dans les quatre cas. Ces
dessins sont des états de référence construits, pas des mesures Monte-Carlo.}

\begin{center}\small
\begin{tabular}{lccc}\toprule
Configuration & $(p_0,p_1,p_2,p_3)$ & $\eta$ & $q$\\\midrule
Au sur le sous-réseau 0 & $(1,0,0,0)$ & $1$ & $1$\\
Au sur le sous-réseau 1 & $(0,1,0,0)$ & $-1/3$ & $1$\\
Au sur le sous-réseau 2 & $(0,0,1,0)$ & $-1/3$ & $1$\\
Au sur le sous-réseau 3 & $(0,0,0,1)$ & $-1/3$ & $1$\\\bottomrule
\end{tabular}\end{center}

Si ces quatre états sont visités aussi souvent, leur occupation moyenne
vaut $(1/4,1/4,1/4,1/4)$. Pourtant, chaque instantané est ordonné :
\[
\boxed{\langle q(p)\rangle=1\quad\text{alors que}\quad
q(\langle p\rangle)=0.}
\]
Les courbes de gauche de la figure suivante sont des moyennes sur les
mesures \emph{et les chaînes}. Leur proximité avec $1/4$ à basse température
reflète la visite de plusieurs variantes ; elle ne suffit pas à conclure
au désordre. Il faut les lire avec $\langle q\rangle$, qui reste proche de 1.

\fig{09_sous_reseaux}{1}{Paramètre signé de l'énoncé, amplitude invariante et
occupations des sous-réseaux. Une moyenne entre variantes ne doit pas être
confondue avec une configuration instantanée désordonnée.}

\begin{oral}{Question 3 -- Une occupation moyenne peut-elle masquer l'ordre ?}
\begin{itemize}
\item \textbf{Partir d'un instantané.} Dans $(1,0,0,0)$, tous les Au
sont sur les sommets ; dans $(0,0,1,0)$, ils occupent un autre sous-réseau.
Les deux configurations ont $q=1$, mais $\eta$ vaut respectivement
$1$ et $-1/3$. Le paramètre du sujet est conservé ; il suit une variante
choisie et ne reconnaît pas les quatre variantes avec le même signe.
\item \textbf{Faire la bonne moyenne.} Si une chaîne ou deux chaînes
indépendantes visitent ces deux états, les occupations moyennes peuvent
valoir $(1/2,0,1/2,0)$. Cela ne signifie pas qu'un instantané contient deux
sous-réseaux à moitié occupés. Il faut moyenner $q$ calculé sur chaque
instantané. La norme des occupations moyennes perd l'information sur
les changements de variante. Une bosse de $p_1(T)$ peut donc provenir
d'un transfert de variante alors que $q$ décroît.
\item \textbf{Tenir compte du volume.} Dans un mélange désordonné fini,
les occupations fluctuent autour de $1/4$ ; leur norme a une moyenne
positive. Au chaud, un $q$ plus faible pour L12 que pour L6 ne signifie
pas que l'algorithme a détruit un ordre stable : une partie du résidu L6
était un effet de taille. La comparaison des trois tailles et des
occupations séparées par chaîne est développée dans le bonus.
\end{itemize}
\end{oral}

\section{Question 4 -- Vérifications du programme}
Le dossier contient une suite de tests automatiques : \TestCount\ tests
ont été relancés avec succès pour la présente mise à jour. Leur date,
leurs valeurs de référence et leur environnement sont précisés ci-dessous.
Ils
contrôlent la géométrie périodique, la réciprocité des voisinages, les quatre
variantes, les valeurs analytiques, les différences d'énergie locales et
globales, la conservation de composition et la reproductibilité.
Les tests supplémentaires vérifient le bootstrap du maximum, sa
reproductibilité et son invariance à une constante ajoutée à l'énergie.
Ils vérifient également la continuité du cycle thermique et l'identité
des trajectoires après reprise à un palier sauvegardé.
Les tests de checkpoints vérifient aussi la reprise exacte au milieu
d'une campagne canonique et le refus d'une configuration incompatible.
Trois tests de l'analyse croisée contrôlent l'identité énergie--Cowley,
la décomposition avec covariance sans ajout de dispersion entre chaînes,
et le comptage des populations sans franchir artificiellement les
frontières entre chaînes.
\input{current_validation.tex}

\subsection{Création d'une paire d'antisites}
Échangeons un Au et un Cu premiers voisins dans l'état $L1_2$. En excluant
leur liaison mutuelle, le premier site a onze voisins Cu, et le second trois
Au et huit Cu. La variation de première coquille vaut
\begin{align*}
\Delta E_1={}&11(J_{\Cu\Cu}-J_{\Au\Cu})
+3(J_{\Au\Au}-J_{\Au\Cu})\\
&+8(J_{\Au\Cu}-J_{\Cu\Cu})=3w_1=0.483\ \mathrm{eV}.
\end{align*}
Douze liaisons identiques de seconde coquille deviennent Au--Cu :
\[
\Delta E_2=12(0.0063)=0.0756\ \mathrm{eV}.
\]
Ainsi
\[
\boxed{\Delta E_{\text{échange 1er voisin}}=0.5586\ \mathrm{eV}.}
\]
L'échange inverse rend exactement l'énergie opposée.

\subsection{Contrôles indépendants}
Le noyau rapide par comptage est confronté à une somme explicite des
matrices de paire pour des échanges aléatoires, voisins et non voisins, sur
plusieurs tailles. Un petit ensemble restreint de configurations est énuméré
exactement pour vérifier la stationnarité canonique de la matrice de
Metropolis. Le signe de l'acceptation des échanges de répliques est également
vérifié par balance détaillée.

Pendant la campagne, à la fin de chaque segment, le code recalcule l'énergie
complète de chaque réplique et impose l'égalité entière avec l'énergie suivie
localement. Il vérifie également la composition et la borne $E\ge E_0$.
Une incohérence arrête immédiatement le calcul.

\begin{oral}{Question 4 -- Qu'est-ce qui valide le code et qu'est-ce qui valide une mesure ?}
\begin{itemize}
\item \textbf{Tester une valeur connue.} Le défaut d'antisites précis
coûte 0.5586 eV. Un test compare cette valeur analytique à deux calculs :
la variation locale rapide et la différence de deux énergies globales.
Un test de géométrie vérifie aussi la réciprocité et les 12 + 6 voisins.
Ces contrôles peuvent déceler une liaison oubliée, un signe erroné ou
un double comptage qu'une courbe d'apparence régulière cacherait.
\item \textbf{Tester la loi et la reprise.} Une petite matrice de
transition énumérée vérifie la mesure canonique dans un espace restreint,
sans résoudre exactement les 864 sites. Les tests de reprise interrompent
puis restaurent un calcul et exigent la même trajectoire que sans arrêt.
Sauvegarder les occupations seules ne suffirait pas : il faut aussi les
générateurs, les énergies, les identités et la parité des échanges.
\item \textbf{Ne pas confondre exactitude et exploration.} Un code correct
peut rester très longtemps dans un bassin. Les tests automatiques et
les contrôles d'archives passent pour L12, alors que ses chaînes
divergent près du pic. R-hat, effectif utile, blocs, demi-séries et
circulation répondent à cette autre question : les mesures disponibles
représentent-elles suffisamment la distribution cible ?
\end{itemize}
\end{oral}

\section{Question 5 -- Énergie de cohésion à température nulle}
Chaque Au a douze voisins Cu. Il y a donc $216\times12=2592$ paires Au--Cu,
aucune paire Au--Au de première coquille, et $5184-2592=2592$ paires Cu--Cu.
Les seconds voisins sont de même espèce et ne contribuent pas à l'énergie.
Le tableau compte les paires \emph{non orientées}, une seule fois chacune :
\begin{center}\small
\begin{tabular}{lrrrr}\toprule
Coquille & Au--Au & Au--Cu & Cu--Cu & Total\\\midrule
Première & 0 & 2592 & 2592 & 5184\\
Deuxième & 648 & 0 & 1944 & 2592\\\bottomrule
\end{tabular}\end{center}
Par exemple, les $216\times6$ contacts Au--Au de seconde coquille sont
comptés depuis les deux extrémités, soit $216\times6/2=648$ paires.
\[
E_0=2592(-0.689)+2592(-0.582)=-3294.432\ \mathrm{eV}.
\]
Donc
\begin{equation}
\begin{aligned}
u_0&=\frac{E_0}{N}=\frac{-3294.432}{864}
=\boxed{-3.813\ \mathrm{eV/atome}},\\
E_0^{\Cu_3\Au}&=\frac{E_0}{N/4}=\frac{-3294.432}{216}=4u_0
=\boxed{-15.252\ \mathrm{eV/formule}}.
\end{aligned}
\end{equation}
Une unité de formule Cu$_3$Au contient trois Cu et un Au, donc quatre
atomes ; la boîte contient $864/4=216$ unités de formule.
La cohésion exprimée comme énergie positive de dissociation vaut
$3.813\ \mathrm{eV/atome}$, relativement au zéro énergétique du modèle.

\subsection{Preuve que \texorpdfstring{$L1_2$}{L12} est un état fondamental}
Soient $A_1,B_1,C_1$ les nombres de paires Au--Au, Au--Cu et Cu--Cu de
première coquille et $M_1=Nz_1/2$ leur nombre total. À composition fixée,
\[
2A_1+B_1=12N_{\Au},\qquad A_1+B_1+C_1=M_1.
\]
En éliminant $B_1$ et $C_1$ dans l'énergie, on trouve
\[
E_1=J_{\Cu\Cu}M_1+(J_{\Au\Cu}-J_{\Cu\Cu})12N_{\Au}+w_1A_1.
\]
Pour les seconds voisins, $E_2=0.0063B_2$. Puisque $L1_2$ possède
$A_1=B_2=0$,
\begin{equation}
\boxed{E-E_0=0.161A_1+0.0063B_2\ge0.}
\end{equation}
C'est une preuve du minimum global pour les paramètres donnés. Elle ne
nécessite aucune simulation.

\begin{oral}{Question 5 -- De la boîte à l'atome puis à la formule}
\begin{itemize}
\item \textbf{Choisir le dénominateur physique.} $E_0$ est l'énergie
de tous les sites ; diviser par $N=864$ donne l'énergie d'un atome en
moyenne, $u_0=-3.813$ eV/atome. Une formule Cu$_3$Au compte quatre atomes :
il existe 216 formules dans la boîte, d'où $E_0/216=4u_0=-15.252$
eV par unité de formule. Pour L12, $E_0$ est huit fois plus grand
en valeur absolue, mais ces quantités normalisées restent les mêmes.
\item \textbf{Lire le signe.} Une énergie de liaison négative abaisse
l'énergie relativement au zéro du modèle. L'énergie positive nécessaire
pour dissocier ce fondamental est son opposé par atome. Ce zéro et les
paramètres de paire appartiennent au modèle ; ils ne remplacent pas une
mesure complète de cohésion de l'alliage réel.
\item \textbf{Distinguer preuve et recherche numérique.} L'identité
$E-E_0=0.161A_1+0.0063B_2$ est positive pour toute configuration de même
composition. Comme $L1_2$ annule les deux termes, il réalise un minimum
global. Observer seulement la plus basse énergie d'une simulation
n'aurait pas constitué cette preuve ; inversement, la preuve du
fondamental ne détermine pas la température de désordre à entropie finie.
\end{itemize}
\end{oral}

\section{Question 6 -- Équilibrer à différentes températures}
La simulation doit échantillonner
\[
P_T(C)=\frac{e^{-\beta E(C)}}{\sum_{C'\,\text{de même composition}}e^{-\beta E(C')}}.
\]
Une configuration finale seule n'est pas une moyenne thermique. À chaque
température, on distingue une phase d'équilibration, écartée de l'analyse,
et une phase de production.

\subsection{Échanges de répliques et mesure canonique}
Pour faciliter les passages entre états métastables, plusieurs répliques
évoluent à des températures $T_r$. La mesure jointe est le produit
$\Pi(\{C_r\})\propto\prod_r e^{-\beta_rE(C_r)}$. L'échange de deux
configurations à températures voisines est accepté avec
\begin{equation}
A_{r,r+1}=\min\left[1,
\exp\{(\beta_r-\beta_{r+1})(E_r-E_{r+1})\}\right].
\end{equation}
Le rapport des poids après et avant échange donne directement cet exposant.
La marginale à chaque température reste canonique. La mesure doit être
classée par \emph{température fixe}, et non par identité de réplique.

Une session d'échange a lieu tous les \ExchangeEvery\ balayages. Le code
alterne les liens pairs et impairs : un lien donné est donc testé toutes
les deux sessions, soit tous les 10 balayages dans cette campagne. Les noyaux de
Metropolis restent ceux demandés par l'énoncé ; les échanges de répliques
ajoutent un mécanisme d'échantillonnage qui conserve la mesure cible
\cite{metropolis,hukushima}.

\subsection{Protocole réellement exécuté}
\begin{center}
\begin{tabular}{lr}
\toprule
Paramètre & Valeur\\
\midrule
Nombre de sites & \Natoms\\
Nombre de températures & \Ntemps\\
Chaînes indépendantes & \Nchains\\
Balayages d'équilibration par chaîne & \BurnSweeps\\
Balayages de production par chaîne & \ProductionSweeps\\
Intervalle entre mesures & \SampleEvery\\
Mesures par température et par chaîne & \Ndraws\\
Maximum sur la grille & $\Tpeak\ \mathrm K$\\
\bottomrule
\end{tabular}
\end{center}
Les graines et toutes les températures figurent dans le fichier de
configuration et dans \texttt{summary.json}. Les chaînes démarrent dans
des variantes ordonnées différentes ; dans la seconde chaîne, les répliques
chaudes démarrent désordonnées. Un départ entièrement aléatoire à très basse
température risquerait de se figer avant d'atteindre l'état fondamental.

\fig{16_coupes_2D}{.95}{Deux plans adjacents des mêmes configurations de la
première chaîne. Les deux hauteurs sont nécessaires pour montrer les quatre
sous-réseaux : une seule coupe peut ne contenir aucun Au dans une variante
pourtant parfaitement ordonnée. L'ordre est plus lisible ici que dans le
nuage tridimensionnel, conservé comme fichier complémentaire
\texttt{07\_configurations.pdf} sans seconde reproduction des mêmes états.}
\fig{17_moyennes_cumulees}{.95}{Comparaison des moyennes cumulées de
production dans les deux chaînes à $\Tpeak$ K. Une moyenne cumulée utilise
toutes les mesures précédentes ; ses points ne sont donc pas indépendants.}

Pour $M'$ mesures de production, la courbe cumulée est
$\overline A(M')=(1/M')\sum_{m=1}^{M'}A_m$. Sa stabilisation et l'accord
des chaînes complètent les traces instantanées, sans prouver à eux seuls
l'absence de bassins métastables non visités. Les contrôles de circulation
donnent 4 puis 3 allers-retours complets froid--chaud--froid pour les deux
chaînes de cette campagne ; leur faible nombre reste une limite à considérer.
\fig{10_repliques}{1}{Trajectoires de quelques identités de répliques.
La circulation en température complète les contrôles locaux d'équilibre.}

\begin{oral}{Question 6 -- Quand peut-on appeler les configurations « d'équilibre » ?}
\begin{itemize}
\item \textbf{Séparer cible et durée.} Metropolis possède la bonne
distribution stationnaire, mais cela ne donne pas le nombre de
balayages nécessaire pour l'approcher. Les états initiaux sont d'abord
relaxés ; leurs mesures sont écartées. La production continue ensuite
sans réinitialisation. Un état final bien ordonné ou une moyenne
stabilisée ne suffisent pas à exclure un bassin jamais visité.
\item \textbf{Expliquer l'aide des répliques chaudes.} Une configuration
peut franchir des barrières à haute température puis rejoindre un
emplacement plus froid par échanges. Les probabilités d'échange
préservent les marginales canoniques. On doit enregistrer
les mesures à température fixe, et suivre séparément les identités
pour mesurer leur circulation. Une même identité n'est pas une série
à température fixe.
\item \textbf{Reconnaître les limites observées.} L6 satisfait les
critères locaux autour du pic, avec seulement sept allers-retours
complets au total. L9 et L12 n'en enregistrent aucun et présentent
des séries critiques insuffisantes. Un taux moyen d'échange élevé
peut masquer un lien limitant. Les résultats sont donc qualifiés
d'estimations canoniques avec diagnostics, plutôt que d'équilibre
garanti par le seul nombre de balayages.
\end{itemize}
\end{oral}

\section{Question 7 -- Grandeurs physiques et incertitudes}
Les moyennes numériques approchent
$\langle A\rangle_T=\sum_C A(C)P_T(C)$. On calcule l'énergie par atome,
les deux paramètres de Cowley, $\eta$, $q$ et les occupations des quatre
sous-réseaux.

\subsection{Capacité calorifique et fluctuations d'ordre}
\input{representative_results.tex}
Ce tableau est calculé à partir du CSV des mesures. La valeur observée
$c_V=0$ au froid signifie qu'aucune fluctuation d'énergie n'a été mesurée
dans cette série ; elle ne démontre pas que la capacité calorifique exacte
est strictement nulle à une température finie.

La capacité calorifique mesure la variation d'énergie moyenne quand on
chauffe à volume constant :
\[
C_V=\left(\frac{\partial\langle E\rangle}{\partial T}\right)_V,
\qquad c_V=C_V/N.
\]
Le déroulement concret est
\[
\text{configurations Metropolis}\ \longrightarrow\
\{E_1,\ldots,E_M\}\ \longrightarrow\
\langle E^2\rangle-\langle E\rangle^2\ \longrightarrow\ c_V.
\]
Metropolis fournit les configurations avec leurs poids thermiques ; le
post-traitement utilise les énergies mesurées, pas le taux d'acceptation.
Dans l'ensemble canonique,
$Z=\sum_C e^{-\beta E(C)}$ et $\langle E\rangle=-\partial_\beta\ln Z$.
En dérivant une seconde fois, on obtient la variance d'énergie.
À Hamiltonien indépendant de la température,
\[
\frac{\partial\langle E\rangle}{\partial\beta}
=-\big(\langle E^2\rangle-\langle E\rangle^2\big).
\]
Comme $d\beta/dT=-1/(\kb T^2)$,
\begin{equation}
\boxed{\frac{c_V}{\kb}=
\frac{\langle E^2\rangle-\langle E\rangle^2}{N(\kb T)^2}.}
\end{equation}
Il faut utiliser l'énergie \emph{totale} dans cette formule. La contribution
calculée est configuratoire : elle n'inclut pas la chaleur des vibrations.

On suit aussi
\[
X_q=\frac{N}{\kb T}\big(\langle q^2\rangle-\langle q\rangle^2\big),
\quad
U_q=1-\frac{\langle q^4\rangle}{3\langle q^2\rangle^2}.
\]
$X_q$ est une réponse à un champ artificiel couplé comme $-Nhq$ ; il ne
faut pas l'identifier sans précaution à une susceptibilité d'un composant
du paramètre d'ordre. Le cumulant est fourni comme diagnostic complémentaire.
Le facteur $3$ est une convention scalaire ; sa limite désordonnée n'est
pas supposée nulle pour cette amplitude à plusieurs composantes.
Par exemple, si les trois composantes de l'ordre sont gaussiennes,
indépendantes et de même variance, leurs moments donnent
$\langle q^4\rangle/\langle q^2\rangle^2=5/3$, donc $U_q=4/9$.
Cette référence explique la valeur chaude proche de 0.44, sans supposer
que toutes les distributions à température finie soient gaussiennes.
$X_q$ est exprimé en eV$^{-1}$ ; $q$, $U_q$ et $c_V/\kb$ sont sans dimension.

\subsection{Autocorrélation, erreurs et comparaison des chaînes}
Les mesures successives ne sont pas indépendantes. Pour une série $A_t$,
\[
\tau_{\mathrm{int}}=\frac12+\sum_{\ell\ge1}\rho_A(\ell),\qquad
N_{\mathrm{eff}}\simeq\frac{M}{2\tau_{\mathrm{int}}}.
\]
L'\textbf{ESS} est ce nombre estimé d'informations indépendantes ;
ce n'est pas le nombre de lignes enregistrées. Le diagnostic
$\widehat R$ compare les distributions des chaînes et de leurs moitiés :
une valeur proche de 1 indique un accord observé, tandis qu'une valeur
élevée signale des moyennes ou des dispersions différentes.
Le \textbf{bootstrap par blocs} rééchantillonne des morceaux contigus
de trajectoire afin de conserver une partie de leur dépendance temporelle.
Une erreur-type, notée SE dans les tableaux, mesure l'incertitude estimée
de la grandeur calculée ; ce n'est pas la fluctuation thermique de cette grandeur.
Le code utilise une troncature par paires positives de l'autocorrélation
et un bootstrap par blocs circulaires. Les blocs ont une longueur d'au moins
$2\tau_{\mathrm{int}}$, estimée à partir de l'énergie et de $q$, sauf lorsque
la série est trop courte pour le permettre ; ce cas est signalé par le faible
nombre de blocs. Les erreurs retiennent aussi la dispersion entre chaînes.
Les variances thermiques sont calculées \emph{dans chaque chaîne} avant
moyennage pour éviter de transformer une divergence entre chaînes en
capacité calorifique artificielle.

Les contrôles indicatifs demandent $\widehat R<1.05$ pour l'énergie et $q$,
un effectif indépendant estimé d'au moins $200$, au moins $20$ blocs par
chaîne, et un écart entre demi-séries inférieur à trois unités d'erreur
estimée. Les graines indépendantes ne rendent pas les températures d'une
même campagne indépendantes.
Ces seuils sont des filtres d'alerte. Deux chaînes restées dans la même
population métastable peuvent présenter un bon accord et manquer ensemble
une autre population. Réussir les contrôles locaux ne certifie donc pas
les poids de toutes les configurations à l'équilibre.

Lorsque E et q sont constants dans une série froide, le code attribue
$\tau=1/2$ et peut obtenir $\widehat R=1$ : ce traitement numérique ne
prouve pas que toutes les excitations ou variantes ont été explorées.
Le choix des blocs utilise E et q ; il reste utile de tester également
la corrélation des paramètres de Cowley et la stabilité des erreurs
lorsque la longueur des blocs est augmentée.

Les paramètres de Cowley sont représentés une seule fois en
figure~\ref{fig:12_cowley}. L'énergie, $q$ et les réponses au voisinage de la
transition sont regroupés en figure~\ref{fig:11_zoom_transition} ; les
occupations et $q$ sur la gamme complète figurent en
figure~\ref{fig:09_sous_reseaux}. Les barres donnent une erreur-type
estimée, pas une enveloppe de l'erreur de modèle. Les exports
\texttt{02\_ordres.pdf} et \texttt{03\_thermodynamique.pdf} restent disponibles,
sans reproduire une deuxième fois ces courbes dans le PDF.
\fig{05_diagnostics}{1}{Diagnostics par température. Un indicateur favorable
ne prouve pas l'exploration de tous les états métastables.}
\fig{06_traces}{1}{Énergie et ordre à température fixe. La phase initiale
grisée est exclue des statistiques.}

\subsection{Vérification de la limite de haute température}
Dans un mélange indépendant de grande taille,
$P(\Au\Au)=1/16$, $P(\Cu\Cu)=9/16$ et $P(\Au\Cu)=3/8$. Ainsi
\begin{align*}
u_\infty={}&6\left(\frac{J_{\Au\Au}}{16}
+\frac{9J_{\Cu\Cu}}{16}+\frac{3J_{\Au\Cu}}8\right)
+3\left(\frac{3I_{\Au\Cu}}8\right)\\
={}&-3.7455375\ \mathrm{eV/atome}.
\end{align*}
À $N$ fini, les probabilités exactes sont
\[
P(\Au\Au)=\frac{N_{\Au}(N_{\Au}-1)}{N(N-1)},\quad
P(\Cu\Cu)=\frac{N_{\Cu}(N_{\Cu}-1)}{N(N-1)},\quad
P(\Au\Cu)=\frac{2N_{\Au}N_{\Cu}}{N(N-1)}.
\]
Le code utilise ces corrections pour les tests. La limite exacte pour la
boîte étudiée vaut $\RandomLimit\ \mathrm{eV/atome}$. Même à plusieurs
milliers de kelvins, les interactions peuvent encore créer des corrélations :
il ne faut pas forcer les courbes à rejoindre la limite $T\to\infty$.

\begin{oral}{Question 7 -- De la trajectoire aux grandeurs et à leurs erreurs}
\begin{itemize}
\item \textbf{Une mesure n'est pas une température.} À T fixé,
Metropolis fournit une succession de configurations. Chaque configuration
donne une énergie totale et des paramètres d'ordre ; on moyenne ces
valeurs pendant la production. $c_V$ demande en plus les fluctuations
de l'énergie : des séjours dans deux populations éloignées produisent
une grande variance, même si leurs énergies moyennes sont stables.
\item \textbf{Utiliser la bonne normalisation.} La formule écrite ici
contient Var$(E)$, puis une division par $N(k_BT)^2$. Si les données
sont déjà en énergie par atome, il faut au contraire multiplier
Var$(u)$ par $N/(k_BT)^2$. Metropolis sert à donner les poids
canoniques ; aucun tirage supplémentaire n'est requis pour calculer
ensuite la variance archivée. Hors équilibre,
cette même variance reste calculable mais n'est pas nécessairement
la dérivée thermique d'équilibre.
\item \textbf{Compter les informations indépendantes.} Des mesures
voisines peuvent décrire presque le même état. Pour $M=3000$ et
$\tau=50$ mesures, l'effectif utile d'une chaîne n'est qu'environ 30,
pas 3000. Les blocs préservent cette dépendance lors du bootstrap.
L'erreur de moyenne et la fluctuation thermique ne sont pas la même
chose : allonger la série peut réduire la première sans faire
disparaître la seconde. Les bassins non visités restent une erreur
systématique que le bootstrap ne reconstruit pas.
\end{itemize}
\end{oral}

\section{Question 8 -- Repérer la transition et comparer au champ moyen}
\subsection{Résultat Monte-Carlo et résolution de la grille}
Le maximum de capacité calorifique sur les températures échantillonnées est
\begin{equation}
\boxed{T^*_{C_V}(L=\Lbox)=\Tpeak\ \mathrm K.}
\end{equation}
Les deux températures voisines de ce maximum sont $\Tleft$ et $\Tright$ K.
L'intervalle $[\Tleft,\Tright]$ K décrit le voisinage sur la grille :
\textbf{ce n'est pas un intervalle de confiance statistique}. Il ne couvre
ni une éventuelle erreur systématique d'équilibration ni le déplacement
avec la taille.

Les maxima trouvés séparément dans les chaînes sont : \ChainPeaks.
Les diagnostics des cinq températures centrées sur le maximum sont résumés
ci-dessous. Leur statut global est : \DiagnosticStatus.
\input{critical_table.tex}
\fig{11_zoom_transition}{1}{Agrandissement de la zone de transition :
ordre, énergie et fluctuations. La zone grisée représente les voisins de
grille, sans interprétation comme intervalle de confiance.}

Une boîte finie n'a pas de singularité thermodynamique. Une seule taille
permet de repérer un maximum apparent, mais pas d'extrapoler $T_c$ à la limite
$N\to\infty$. La comparaison des trois tailles en
section~\ref{sec:tailles} examine distributions et déplacements des maxima,
avec les réserves d'échantillonnage nécessaires \cite{binder}.

\fig{04_distributions}{1}{Distributions près de la zone de transition.
L'inspection de la coexistence complète celle du maximum de capacité
calorifique ; les histogrammes restent sensibles à la durée du calcul.}

\subsection{Incertitude statistique du maximum et concordance des indices}
La hauteur du pic possède une erreur-type, mais celle-ci n'est pas une
erreur en kelvins sur sa position. On estime cette dernière par un
bootstrap \emph{conjoint} des températures : dans chaque chaîne, les mêmes
indices de blocs sont appliqués simultanément à toutes les températures.
Les deux chaînes sont rééchantillonnées séparément. Pour chaque tirage,
on recalcule les variances d'énergie dans chaque chaîne, on moyenne leurs
capacités calorifiques puis on recherche le maximum sur la grille initiale.
Cette synchronisation conserve les dépendances entre températures induites
par les échanges de répliques.

La longueur de référence vaut
$b=\lceil2\max_{T,c,A\in\{E,q\}}\tau_{c,A}\rceil$, avec les bornes
numériques décrites en annexe. On la compare à une longueur environ deux
fois plus petite et deux fois plus grande. Les résultats sont obtenus
uniquement à partir des archives de production :
\input{peak_uncertainty.tex}
\fig{18_incertitude_pic}{.95}{Fréquence de la position du maximum parmi
les tirages bootstrap pour trois longueurs de bloc. Les températures
possibles sont celles de la grille : aucune précision sous la maille de
température n'est créée par ce calcul.}

Un écart-type inférieur au pas de 5 K traduit ici des changements rares
de la case de maximum ; il ne localise pas le pic entre les températures
simulées. L'intervalle percentile est une estimation statistique conditionnelle à
la taille, aux deux chaînes observées et à la grille. Il ne couvre ni
le déplacement avec $L$, ni les états métastables éventuellement non
visités. S'il se réduit à un point, cela ne signifie pas que la vraie
température de transition est connue exactement. Les changements de pente
et les maxima de fluctuations ne définissent pas exactement le même
critère : leur proximité identifie une zone cohérente, pas plusieurs
mesures indépendantes d'une température critique unique.

Une étude de plusieurs tailles à précision statistique contrôlée
renforcerait l'interprétation thermodynamique. Les résultats de cette
question concernent $L=6$.
Les campagnes $L=9$ et $L=12$ sont achevées et techniquement validées ;
leurs résultats, avec leurs réserves statistiques, sont présentés en
section~\ref{sec:tailles}.

\subsection{Estimation analytique de Bragg-Williams}
Pour une famille d'états à un sous-réseau privilégié,
avec $0\leq\eta\leq1$ dans cette comparaison d'amplitudes,
\[
p_0=\frac{1+3\eta}{4},\qquad
p_1=p_2=p_3=p_f=\frac{1-\eta}{4}.
\]
Le champ moyen néglige les corrélations entre occupations. Les premiers
voisins relient des sous-réseaux distincts ; les seconds voisins appartiennent
au même sous-réseau. Le comptage donne
\[
\frac{B_1}{N}=\frac94+\frac34\eta^2,\qquad
\frac{B_2}{N}=\frac98(1-\eta^2).
\]
L'énergie moyenne s'écrit
\[
u_{\mathrm{MF}}(\eta)=u_\infty-A\eta^2,
\quad A=\frac38w_1-\frac9{16}w_2=0.0674625\ \mathrm{eV},
\]
avec $w_2=I_{\Au\Au}+I_{\Cu\Cu}-2I_{\Au\Cu}=-0.0126$ eV.

L'entropie de mélange par site d'un sous-réseau vaut $-\kb h(p)$, où
$h(p)=p\ln p+(1-p)\ln(1-p)$ et $0\ln0=0$. Donc
\[
f_{\mathrm{MF}}(\eta,T)=u_\infty-A\eta^2+\kb T g(\eta),
\quad g(\eta)=\tfrac14h(p_0)+\tfrac34h(p_f).
\]
La coexistence vérifie à la fois la stationnarité du minimum ordonné et
l'égalité de son énergie libre à celle du minimum désordonné :
\[
2A\eta_t=\kb T_tg'(\eta_t),\qquad
A\eta_t^2=\kb T_t[g(\eta_t)-g(0)].
\]
En éliminant $T_t$,
\[
\eta_tg'(\eta_t)=2[g(\eta_t)-g(0)],\qquad
g'(\eta)=\frac3{16}\ln\frac{(1+3\eta)(3+\eta)}{3(1-\eta)^2}.
\]
La résolution donne $\eta_t\simeq0.46358$ et
\[
\boxed{T_t^{\mathrm{MF}}\simeq1712.18\ \mathrm K.}
\]
L'annulation de $f''(0,T)$ donne seulement la perte de stabilité locale de
l'état désordonné : $T_{\mathrm{sp}}=2A/\kb\simeq1565.74$ K. Cette
spinodale est distincte de la coexistence. Le saut d'ordre prédit ici est
une propriété de l'approximation de champ moyen, et ne remplace pas une
étude de la nature de la transition du modèle spatial.

\fig{08_champ_moyen}{1}{Énergies libres approchées et comparaison entre
amplitudes d'ordre. L'estimation de champ moyen n'est pas utilisée pour
fixer le résultat Monte-Carlo.}

\begin{oral}{Question 8 -- Quelle température avons-nous effectivement déterminée ?}
\begin{itemize}
\item \textbf{Nommer le critère.} Le maximum de $c_V$ sur la grille
définit $T^*(L)$, un repère de boîte finie. Les voisins de grille
décrivent la résolution, tandis que le bootstrap décrit les déplacements
de la case de maximum compatibles avec les trajectoires observées.
Ni l'un ni l'autre ne couvre un bassin absent, une extrapolation en
volume ou l'erreur du Hamiltonien. L12 donne ainsi 1052.5 K observés,
avec une statistique encore insuffisante au maximum.
\item \textbf{Croiser les symptômes.} Une chute de $q$, une montée de
l'énergie, des fluctuations fortes et deux populations sont cohérentes
avec une transition dans la même zone. Elles utilisent toutefois
les mêmes mesures, et énergie et Cowley sont liés exactement.
Les maxima de $X_q$ ou les creux du cumulant ne sont donc pas
des expériences indépendantes donnant chacune une valeur exacte de $T_c$.
\item \textbf{Comprendre le champ moyen.} Bragg--Williams remplace
les voisinages corrélés par des probabilités indépendantes. Il peut
donner une coexistence, où deux minima d'énergie libre sont égaux,
et une spinodale, où l'un perd sa stabilité locale. Ces températures
diffèrent entre elles et du maximum Monte-Carlo fini. L'écart avec
Monte-Carlo renseigne sur l'approximation ; il ne suffit pas à
accuser l'implémentation. Aucune de ces températures ne devient
automatiquement celle de Cu$_3$Au réel.
\end{itemize}
\end{oral}

\section{Question 9 -- Discussion des résultats et limites}
\input{discussion_results.tex}

Le mécanisme général est une compétition entre énergie et entropie :
$F=U-TS$. L'ordre minimise l'énergie, tandis que le nombre très élevé de
configurations désordonnées favorise leur poids statistique à haute
température. La disparition de l'ordre global peut précéder la disparition
des préférences chimiques locales : un $q$ faible ne contraint pas
$\alpha_1$ et $\alpha_2$ à être exactement nuls.

Les limites principales sont explicites : tailles finies, durées de
simulation finies, grille de températures discrète, et approximation physique
du réseau rigide. Les vibrations, la dilatation et les relaxations autour des
antisites sont absentes. Les constantes de paire sont supposées indépendantes
de la température. On ne présente donc pas la température du modèle comme
une prédiction expérimentale exacte.

\subsection{Comparaison critique avec \texorpdfstring{Cu$_3$Au}{Cu3Au} réel}
Les mesures de calorimétrie et de diffusion situent le désordre chimique de
Cu$_3$Au près de 660--663 K, et non à 1050 K
\cite{benisek,paras}. L'écart d'environ 390 K est de l'ordre de 60\,\% :
le maximum du présent modèle est donc qualitativement pertinent mais
quantitativement non calibré au matériau. Passer de $L=6$ à $L=9$ ou $12$ ne
constitue pas une correction des interactions du Hamiltonien. On ne peut
pas attribuer cet écart à la seule taille finie, ni aux seules vibrations,
sans calculs de contrôle séparant ces contributions.
Un maximum calorimétrique expérimental dépend aussi du protocole :
Benisek et Dachs rapportent un pic à 675 K pour leur échantillon le
plus ordonné et une réaction vers 660 K pour des échantillons initialement
désordonnés \cite{benisek}. Il faut donc distinguer température de
transition de référence et maximum mesuré pendant une histoire thermique.
La transition réelle est décrite comme du premier ordre
\cite{paras}. L'arrondi des courbes dans notre boîte finie ne démontre
donc pas une transition continue, et l'on ne doit pas exiger une divergence
critique de $\xi$ pour conclure à l'existence d'une transition.

\begin{center}\small
\begin{tabular}{>{\raggedright\arraybackslash}p{.22\linewidth}>{\raggedright\arraybackslash}p{.34\linewidth}>{\raggedright\arraybackslash}p{.33\linewidth}}
\toprule
Observation ou attente physique & Ce que donne le modèle & Limite ou interprétation\\\midrule
Ordre chimique au froid & $L1_2$, quatre variantes, défauts d'antisites et
préférence pour les liaisons Au--Cu & Accord structurel et énergétique
pour ce modèle ; les positions atomiques ne se relaxent pas.\\
Ordre local après perte d'ordre global & $q$ devient faible alors que
$\alpha_1,\alpha_2$ restent mesurables & Accord qualitatif avec la diffusion
expérimentale sensible à l'ordre à courte portée \cite{owen}.\\
Chaleur mesurée en calorimétrie & Pic de la seule contribution
configuratoire $c_V$ & Une expérience mesure généralement $C_P$, avec
vibrations et autres contributions ; ni le fond thermique, ni le pic
total ne sont prédits ici \cite{benisek}.\\
Diffusion et croissance de domaines & Mise en ordre rapide avec des
échanges Au--Cu à toute distance & Pas de lacunes ni de barrières de saut
calibrées ; aucun temps physique ni coefficient de diffusion ne peut
être déduit du GIF \cite{kessler}.\\
Volume, surfaces et interfaces & Boîte périodique, volume fixé et interfaces
chimiques possibles & Pas de surface libre, joint de grain ou contrainte
élastique relaxable ; les images périodiques peuvent interagir.\\
Très haute température & Le réseau reste artificiellement solide même
à 4000 K & Cette borne sert à explorer la limite mathématique désordonnée ;
le modèle n'a aucun mécanisme de fusion et ne représente pas un solide réel à cette température.\\
\bottomrule
\end{tabular}\end{center}

\textbf{Ce qui est observé malgré les simplifications.} La compétition
énergie--entropie, les quatre variantes, la création de défauts, la perte
d'ordre global et la persistance de préférences locales sont des mécanismes
capturés par les occupations et les interactions de paires. Leur présence
ne valide pas automatiquement leurs amplitudes, leur température ou
leur vitesse dans le matériau réel.

\textbf{Ce que les animations ne résolvent pas.} Selon la température de
trempe et la dynamique choisie, un modèle d'alliage peut présenter une
incubation par nucléation ou une mise en ordre plus homogène
\cite{kessler}. Une seule trajectoire, sans temps d'attente apparent à
la résolution de 20 balayages, ne permet pas de choisir entre ces scénarios.
Les échanges globaux peuvent faciliter une réorganisation collective que
la diffusion locale réaliserait autrement. Inversement, un pic ou une
bosse d'occupation de sous-réseau peut provenir d'un changement de variante,
sans constituer un germe géométrique nouveau.

\subsection{Approximations statistiques et spatiales à ne pas confondre}
Un $q$ résiduel au chaud n'est pas nécessairement un défaut : $q$ est une
norme positive, dont la moyenne ne s'annule pas exactement dans une petite
boîte désordonnée. Une absence de fluctuations au froid borne seulement
ce qui a été vu pendant cette production. Les erreurs par blocs corrigent
la dépendance temporelle observée ; elles ne reconstruisent pas un bassin
métastable jamais visité.

Les longueurs du bonus sont celles d'une enveloppe d'ordre définie par
maille, avec moyenne sur les trois composantes et sur les directions
équivalentes. Ce choix fournit un indicateur scalaire ; il peut masquer
une anisotropie des domaines. Le premier vecteur accessible à $L=6$ vaut
$k_{\min}a\simeq1.047$, loin d'une résolution continue de très petits $k$.
Les distances périodiques et la petite gamme de $r$ interdisent un ajustement
robuste d'une décroissance exponentielle ou d'un exposant critique.
Une longueur canonique supérieure à la boîte peut refléter le mélange
des variantes et l'ordre uniforme ; ce n'est pas la taille d'un domaine
mesuré dans les coupes. La longueur intra-variante est un diagnostic
conditionnel, sans intervalle de confiance calculé dans cette étude.

\subsection{Ce que les tailles supplémentaires permettent de tester}
Le côté passe de $6a$ à $9a$, soit un facteur 1.5 ; le nombre de sites
passe de 864 à 2916, soit un facteur 3.375, avec 729 Au et 2187 Cu.
Le plus petit décalage réciproque passe de $1.047/a$ à $0.698/a$.
Cette augmentation peut réduire l'arrondi des pics et l'influence des
images périodiques, et offrir plus de place aux domaines. Elle ne garantit
aucune amélioration de l'équilibre : les temps de relaxation et les
barrières peuvent augmenter, tandis que les échanges de répliques
deviennent moins efficaces à grille inchangée.

La comparaison L6/L9 conserve le Hamiltonien, les 46 températures,
6000 balayages initiaux et 30000 de production par chaîne. Les moyennes
sont comparées par atome ; les maxima $T^*(L)$, leurs intervalles,
les histogrammes, R-hat, ESS et la circulation des répliques sont
examinés ensemble en section~\ref{sec:tailles}. Les valeurs de $q$ ne sont pas comparées comme
si leur biais de norme finie était indépendant de $N$. Des états
atomiques espacés de 100 balayages ont été conservés pour permettre
une analyse spatiale après le calcul. Ce budget identique de balayages
n'est pas une garantie de nombres identiques d'états indépendants.

L12 porte ensuite la boîte à 6912 sites, avec un plus petit vecteur
accessible $k_{\min}a=0.524$. Sa grille est densifiée et sa production
doublée ; les effets de taille et de protocole ne sont donc pas séparés
parfaitement. Les trois tailles permettent d'observer des tendances,
mais leurs diagnostics critiques inégaux ne justifient ni extrapolation
quantitative ni loi universelle. L'analyse croisée de la
section~\ref{sec:tailles} détaille ce que l'augmentation améliore et
les défauts d'échantillonnage qu'elle n'a pas supprimés.
L'absence de phonons, de dilatation, de lacunes et de dynamique de diffusion
reste la même pour les trois tailles.

Une distribution à deux pics est compatible avec une coexistence ; une
hystérésis est compatible avec une métastabilité. Aucune de ces observations
isolées n'est une preuve suffisante de l'ordre thermodynamique de la
transition. Il faudrait renforcer les statistiques et comparer plusieurs
tailles pour conclure sur ce point.

\begin{samepage}
\subsection{Conclusion physique}
Au froid, l'énergie rejoint le minimum exact
$u_0=-3.813$ eV/atome, $q\simeq1$ et les deux paramètres de Cowley
prennent les valeurs de $L1_2$. En chauffant, les défauts chimiques
augmentent et l'ordre global chute dans la zone du maximum de capacité
calorifique, situé à $\Tpeak$ K sur cette grille. Au chaud, $q$ devient
faible alors que les corrélations locales restent mesurables : perdre
l'ordre à longue portée ne rend pas immédiatement les occupations
indépendantes. La valeur $\Tpeak$ K caractérise donc ce modèle de réseau
rigide dans la boîte de $\Natoms$ sites ; elle n'est ni une extrapolation
à taille infinie ni une mesure expérimentale de Cu$_3$Au.
Les grandes tailles situent elles aussi leur maximum dans la zone
1050--1055 K, tout en révélant un ralentissement important. Leur ajout
consolide les tendances qualitatives ; il ne supprime pas l'incertitude
sur les poids des populations ni l'erreur de modèle.
\end{samepage}

\begin{oral}{Question 9 -- Que retrouve-t-on du matériau réel et que manque-t-il ?}
\begin{itemize}
\item \textbf{Les accords qualitatifs ont une origine identifiable.}
Les liaisons de première coquille favorisent les voisins chimiques
différents et celles de seconde coquille favorisent les espèces identiques.
Le modèle peut donc retrouver la structure $L1_2$, les quatre variantes,
les antisites, une perte d'ordre global et un ordre local persistant.
Ces résultats expliquent le mécanisme énergie--entropie sans supposer
que toutes leurs amplitudes sont quantitativement correctes.
\item \textbf{Une plus grande boîte ne change pas la physique absente.}
Dans le matériau réel, les atomes vibrent et les positions se relaxent,
le volume varie, les défauts migrent par des mécanismes locaux et les
contributions électroniques peuvent compter. Ici les positions sont
fixes, les paires sont constantes et les échanges peuvent relier deux
sites éloignés. L9 puis L12 améliorent l'espace disponible, pas ces
hypothèses. Le décalage entre 660--663 K expérimental et environ
1050--1055 K dans le modèle ne peut donc pas être corrigé par la seule
augmentation du nombre d'atomes.
\item \textbf{La critique concerne aussi les résultats numériques.}
L6 fournit le repère le mieux contrôlé statistiquement près du maximum ;
les grandes boîtes montrent un pic plus étroit et des populations plus
distinctes, mais elles explorent moins bien leur coexistence. La conclusion
doit garder les deux constats : la tendance de taille est intéressante,
son extrapolation n'est pas validée. De même, le GIF montre une mise en
ordre sous la dynamique choisie ; il ne mesure pas une vitesse de
nucléation dans l'alliage réel.
\end{itemize}
\end{oral}

\begin{oral}{Fil conducteur pour une présentation de 8 à 10 minutes}
\begin{enumerate}
\item Décrire le réseau fixe, la composition et la structure $L1_2$.
\item Montrer les préférences de paire et la preuve du fondamental.
\item Expliquer un échange et la balance détaillée de Metropolis.
\item Définir Cowley et l'ordre global ; montrer les quatre variantes.
\item Présenter les tests avant les résultats numériques.
\item Montrer énergie, capacité calorifique et perte d'ordre.
\item Donner le maximum apparent avec sa résolution et ses diagnostics.
\item Terminer par les limites physiques et les diagnostics des trois tailles.
\end{enumerate}
\end{oral}

\section{Bonus -- Histoire thermique, visualisation et corrélations spatiales}
\label{sec:bonus-cycle}
Le bonus étudie la mémoire du trajet thermique à budget de simulation
fixé. Il complète l'étude canonique précédente par un protocole
séquentiel de chauffage puis de refroidissement.
\subsection{Partie 1 -- Cycle thermique et visualisation des configurations}
\begingroup\sloppy
\IfFileExists{bonus_cycle.tex}{\input{bonus_cycle.tex}}{
Le module \texttt{run\_thermal\_cycle.py} produit les mesures et les
comparaisons des deux branches ; aucun résultat de cycle n'est joint
à cette version du rapport.}
\IfFileExists{bonus_animations.tex}{\input{bonus_animations.tex}}{}
\begin{oral}{Bonus 1 -- Comment lire un cycle et une animation ?}
\begin{itemize}
\item \textbf{Lire le temps dans le bon sens.} Au chauffage, les
températures sont parcourues de gauche à droite ; au refroidissement,
le temps avance lorsque T diminue. Le GIF thermique montre des fins
de paliers, séparées par beaucoup de balayages, et ne résout pas les
événements entre ces images. Le GIF isotherme conserve T à 1000 K :
seul le temps Monte-Carlo avance. Les neuf coupes sont des vues du même
état, pas neuf trajectoires indépendantes.
\item \textbf{Relier image, ordre et histoire.} Une couche pauvre en Au
peut appartenir à un cristal ordonné dont les Au occupent les couches
voisines. On confronte donc les coupes à $q$ et aux occupations, séparées
par cycle. Une différence chauffage--refroidissement montre une mémoire
à budget fini. Elle doit être vérifiée en allongeant les paliers, et
ne constitue ni une nouvelle température d'équilibre ni une preuve
isolée de nucléation. Seules des branches réellement calculées peuvent
être comparées ; les productions canoniques L9/L12 n'en sont pas.
\end{itemize}
\end{oral}
\subsection{Partie 2 -- Longueur de corrélation et analyse de Fourier}
\IfFileExists{bonus_correlation.tex}{\input{bonus_correlation.tex}}{
Le module de corrélations produit les champs d'ordre locaux, les spectres
et les diagnostics de la longueur de second moment.}
\begin{oral}{Bonus 2 -- Que mesure la longueur de corrélation ?}
\begin{itemize}
\item \textbf{Poser une question spatiale.} Deux configurations de même
énergie ou de même Cowley peuvent répartir leur ordre sur des distances
différentes. Le champ signé par maille distingue les quatre variantes ;
sa corrélation demande si connaître l'organisation ici renseigne encore
sur celle d'une maille éloignée. Le mode zéro utilisé est celui de
l'ordre chimique, pas celui du nombre total d'Au, fixé par construction.
\item \textbf{Comprendre le rapport de Fourier.} Un grand pic au centre
comparé à ses voisins réciproques correspond à un ordre cohérent sur
de grandes distances. Le second moment convertit ce rapport en une
longueur, avec une convention de réseau explicite. Dans une petite
boîte, un ordre uniforme ou des changements collectifs de variante
peuvent rendre cette longueur supérieure au côté : ce n'est pas un
rayon géométrique. Les fluctuations autour d'une variante, la corrélation
canonique et la croissance hors équilibre sont trois objets distincts.
\item \textbf{Lire ensemble la valeur et sa validité.} Une racine non
réelle reste manquante ; une grande valeur n'est pas écrêtée pour la
rendre plausible. On examine les blocs, les composantes signées et
$\xi/(La)$ avant une interprétation. Les trois tailles améliorent la
résolution des modes accessibles, mais leurs budgets et leurs mélanges
de variantes différents interdisent ici un exposant critique fiable.
\end{itemize}
\end{oral}
\subsection{Partie 3 -- Analyse croisée et superpositions L6/L9/L12}
\label{sec:tailles}
\IfFileExists{bonus_sizes.tex}{\input{bonus_sizes.tex}}{
Les résultats de taille supplémentaire ne sont pas disponibles dans cette version.}
\IfFileExists{bonus_size_cycles.tex}{\input{bonus_size_cycles.tex}}{}
\endgroup

\appendix
\section{Correspondance entre les questions et les fichiers}
\begin{center}\small
\begin{tabular}{p{.12\linewidth}p{.40\linewidth}p{.38\linewidth}}
\toprule
Question & Réponse numérique & Fichier principal\\
\midrule
1 & Réseau, énergie, échanges, Metropolis & \texttt{lattice.py}, \texttt{model.py}, \texttt{kernel.py}\\
2 & Cowley des deux coquilles & \texttt{model.py : cowley}\\
3 & Ordre signé, amplitude, sous-réseaux & \texttt{model.py : ordre}\\
4 & Tests analytiques et probabilistes & \texttt{tests/test\_physics.py}\\
5 & Énergie exacte, preuve et tests & \texttt{model.py}, présent rapport\\
6 & Configurations, répliques, graines & \texttt{simulation.py}, \texttt{chain\_*.npz}\\
7 & Moyennes, erreurs, figures & \texttt{statistics.py}, \texttt{figures.py}\\
8 & Maximum sur la grille, champ moyen & \texttt{summary.json}, \texttt{meanfield.py}\\
9 & Discussion fondée sur les données & Rapport et diagnostics CSV\\
10 & Cycle chauffage--refroidissement & \texttt{thermal\_cycle.py}, \texttt{run\_thermal\_cycle.py}\\
\bottomrule
\end{tabular}\end{center}

Le lanceur \texttt{run\_project.py} enchaîne les tests, la simulation, les
statistiques, les figures et le rapport. L'option \texttt{--reuse} réutilise
uniquement une campagne dont la configuration correspond exactement.
La compilation PDF demande une installation LaTeX, mais les fichiers
\texttt{.tex} et les figures sont produits même sans compilation.

\clearpage
\section{Algorithme}
\begingroup\small
\label{sec:algorithme}

Cette partie décrit le calcul effectivement implémenté, depuis la construction
du réseau jusqu'aux erreurs statistiques. Les algorithmes numérotés forment un
pseudocode commenté : les noms de fonctions permettent de retrouver chaque
opération dans les sources. Les commentaires expliquent à la fois le choix
physique et le moment où l'opération doit être exécutée.

\subsection{Notations, données et niveaux de boucle}
On note $C$ le nombre de chaînes indépendantes, $R$ le nombre de températures,
$S_{\mathrm{eq}}$ et $S_{\mathrm{prod}}$ les nombres de balayages des deux
phases, $d_m$ l'intervalle de mesure, $d_x$ l'intervalle d'échange et $K$ la
taille maximale d'un segment. Dans une chaîne, les $R$ configurations sont
couplées par les échanges de répliques ; ce ne sont donc pas $R$ chaînes
indépendantes. L'indice $r=0,\ldots,R-1$ désigne un \emph{emplacement à
température fixe} $T_r$, avec $T_0<\cdots<T_{R-1}$.

\begin{center}\small
\begin{tabular}{p{.27\linewidth}p{.64\linewidth}}
\toprule
Donnée en mémoire & Signification et unité\\
\midrule
\texttt{states[r,i]} & Occupation $s_{r,i}\in\{0,1\}$ : Cu ou Au.\\
\texttt{neighbors[i,k]} & Indices des 12 premiers puis des 6 seconds voisins.\\
\texttt{sub[i]} & Sous-réseau $\mu\in\{0,1,2,3\}$ du site $i$.\\
\texttt{energies[r]} & Énergie totale entière $\mathcal E_r$ en micro-eV.\\
\texttt{labels[r]} & Identité de la configuration présente à $T_r$.\\
\texttt{phase} & Nombre de sessions d'échange déjà exécutées ; sa parité choisit les liens.\\
\texttt{records[m,r,:]} & Neuf observables à la mesure $m$ et à la température $T_r$.\\
\bottomrule
\end{tabular}
\end{center}

Posons $Q=10^6$, le nombre de micro-eV par eV. Ainsi
$E_r=\mathcal E_r/Q$ en eV et
\[
 \widetilde\beta_r=\frac{1}{Q\kb T_r},\qquad
 \widetilde\beta_r\Delta\mathcal E=\frac{\Delta E}{\kb T_r}.
\]
Cette conversion est indispensable : utiliser une énergie entière en
micro-eV directement avec $1/(\kb T)$ fausserait l'exposant d'un facteur $10^6$.
Le code nomme \texttt{beta} la quantité $\widetilde\beta$.

L'emboîtement exact est
\[
\boxed{\text{chaîne}\ \to\ \text{phase}\ \to\ \text{segment}\ \to\
\text{balayage}\ \to\ \text{température}\ \to\ \text{proposition}.}
\]
À la fin de chaque balayage viennent, dans cet ordre, les échanges de
répliques puis les mesures si leurs échéances sont atteintes. Les statistiques
de production sont calculées après les simulations, à partir des séries
sauvegardées, et non après chaque proposition.

\subsection{Algorithme 1 : préparer la géométrie et les configurations}
\begin{algocode}{Algorithme 1 --- Initialisation d'une chaîne}
\textbf{Entrées :} configuration JSON et indice de chaîne $c=0,\ldots,C-1$.
\textbf{Sorties :} géométrie fixe, occupations, énergies et identités initiales.
\begin{enumerate}
\item Vérifier $L\ge3$, les températures positives strictement croissantes,
au moins deux graines distinctes et les longueurs de phases strictement
positives. Exiger que $S_{\mathrm{eq}}$, $S_{\mathrm{prod}}$ et $K$ soient
divisibles par $d_m$ et, si $d_x>0$, par $d_x$.
\par\textit{Commentaire :} ces multiples empêchent le redémarrage du compteur
local de balayages à chaque segment de décaler les mesures ou les échanges.
Le dernier segment, éventuellement plus court, conserve cette propriété.
\item Construire une seule fois les $N=4L^3$ coordonnées entières
$\mathbf X=2(i,j,k)+\mathbf b_\mu$, avec
$\mathbf b_\mu\in\{(0,0,0),(0,1,1),(1,0,1),(1,1,0)\}$.
Attribuer \texttt{sub[i]} et construire la correspondance coordonnées--indice.
\par\textit{Commentaire :} l'arithmétique entière évite les tests d'égalité
approchée sur des coordonnées. La géométrie est commune à toutes les chaînes.
\item Pour chaque site, ajouter les déplacements des deux coquilles,
réduire chaque coordonnée modulo $2L$ et stocker les 18 indices voisins.
\par\textit{Commentaire :} la périodicité supprime les surfaces libres.
Les listes restent valables pendant tout le calcul puisque les atomes
échangent leurs espèces sans déplacer les sites.
\item Initialiser un générateur avec la graine de la chaîne $c$.
Pour chaque température $T_r$, choisir l'état initial suivant le code :
\begin{itemize}
\item si $c=0$ ou $T_r<T_{\text{départ}}$, remplir d'Au le sous-réseau
$\mu=(2c)\bmod4$ et de Cu les trois autres ;
\item sinon, tirer sans remise exactement $N/4$ sites Au parmi les $N$ sites.
\end{itemize}
\textit{Commentaire :} les deux chaînes de production commencent dans les
variantes 0 et 2. Les départs aléatoires de la seconde chaîne concernent
$T_r\ge1060$ K. Les départs froids ordonnés évitent de partir d'un état
aléatoire difficile à réorganiser à basse température ; ils ne dispensent
pas de vérifier l'équilibration.
\item Calculer pour chaque configuration l'énergie entière de référence
\[
\mathcal E_r=\frac12\sum_{i=0}^{N-1}
 \left[\sum_{k=0}^{11}J^{\mathrm{int}}_{s_{r,i},s_{r,V(i,k)}}
 +\sum_{k=12}^{17}I^{\mathrm{int}}_{s_{r,i},s_{r,V(i,k)}}\right],
\]
où $J^{\mathrm{int}}=QJ$ et $I^{\mathrm{int}}=QI$.
Initialiser \texttt{labels[r]} $=r$ et \texttt{phase} $=0$.
\par\textit{Commentaire :} le facteur $1/2$ corrige le double comptage des
paires orientées. Ce calcul global fournit l'énergie initiale ; les
propositions ultérieures utilisent seulement une différence locale.
\end{enumerate}
{\raggedright\textbf{Sources :} \texttt{simulation.py : load\_config, campaign} ;
\texttt{lattice.py} ; \texttt{model.py : energy\_int}.\par}
\end{algocode}

\subsection{Algorithme 2 : organiser équilibration et production}
\begin{algocode}{Algorithme 2 --- Campagne complète et contrôles par segment}
\textbf{Entrées :} géométrie, paramètres et graines.
\textbf{Sorties :} archives par chaîne, moyennes, erreurs et diagnostics.
\begin{enumerate}
\item \textbf{Pour chaque chaîne} $c=0,\ldots,C-1$ : exécuter l'algorithme 1
et créer des listes séparées pour les mesures d'équilibration et de production.
\par\textit{Commentaire :} chaque chaîne possède sa graine et son propre
ensemble de $R$ répliques ; aucun échange n'a lieu entre chaînes.
\item \textbf{Pour chaque phase}, d'abord équilibration de longueur
$S_{\mathrm{eq}}$, puis production de longueur $S_{\mathrm{prod}}$ :
\begin{enumerate}
\item Mettre le compteur $\texttt{done}=0$.
\item \textbf{Tant que} \texttt{done} est inférieur à la longueur de phase :
choisir $S=\min(K,\text{longueur}-\texttt{done})$, tirer une graine de
segment dans $\{1,\ldots,2^{31}-2\}$ avec le générateur de la chaîne,
puis appeler l'algorithme 3 pour $S$ balayages.
\par\textit{Commentaire :} le noyau Numba réinitialise son générateur à chaque
appel avec cette graine nouvellement tirée. On conserve en revanche
\texttt{states}, \texttt{energies}, \texttt{labels} et \texttt{phase} entre
appels ; les segments découpent une même trajectoire, sans recréer les états.
\item Après chaque segment, recalculer indépendamment toutes les énergies
avec les matrices $J^{\mathrm{int}}$ et $I^{\mathrm{int}}$. Arrêter si elles
diffèrent des énergies mises à jour, si $\sum_i s_{r,i}\ne N/4$, ou si
$\mathcal E_r<QN(-3.813)$ pour une réplique.
\par\textit{Commentaire :} ces contrôles exacts détectent une erreur de
variation locale, une mauvaise permutation ou une violation de composition.
Les effectuer par segment évite un recalcul global à chaque tentative.
\item En équilibration, ajouter les mesures à \texttt{burn\_samples} seulement.
En production, ajouter mesures et identités aux séries d'analyse, et cumuler
les compteurs d'acceptation locale et d'échange. Augmenter \texttt{done} de $S$.
\par\textit{Commentaire :} les mesures initiales servent aux historiques
mais n'entrent pas dans les moyennes thermiques. Le passage à la production
continue avec les configurations atteintes, sans nouvelle initialisation.
\end{enumerate}
\item Concaténer les segments et sauvegarder \texttt{chain\_XX.npz} :
mesures de production, mesures d'équilibration, identités, configurations
finales, températures, coordonnées et sous-réseaux.
\item Après toutes les chaînes, analyser les séries avec les algorithmes 7
et 8 ; écrire \texttt{observables.csv} et \texttt{summary.json}.
\par\textit{Commentaire :} conserver les données brutes permet de refaire
les analyses et les figures sans relancer les trajectoires Monte-Carlo.
\end{enumerate}
\textbf{Source :} \texttt{simulation.py : campaign}.
\end{algocode}

\subsubsection*{Captures spatiales et reprise de la campagne L12}
Les options supplémentaires de L9 et L12 n'altèrent pas la règle
de Metropolis. Une capture recopie les occupations après la mesure au
balayage choisi ; ses dates permettent de retrouver exactement les
observables correspondantes. À L12, les captures tous les 200 balayages
fixent aussi les limites de segments de production, dont la longueur
maximale reste une limite et non une durée imposée.

La reprise optionnelle sauvegarde après chaque segment validé les états,
les énergies, les identités, la parité, la position dans les phases,
les mesures nouvelles et l'état du générateur hôte. Le noyau reçoit
une nouvelle graine à chaque segment : restaurer le générateur hôte
et les limites de segment restaure donc aussi la suite des graines.
Chaque checkpoint est écrit dans un fichier temporaire puis publié
atomiquement ; une écriture inachevée n'est pas utilisée. Une signature
des paramètres, des sources et de l'environnement empêche de reprendre
silencieusement avec un autre noyau. Les tests comparent la trajectoire
reprise à la trajectoire ininterrompue. Le pseudocode P2 décrit le chemin
sans reprise ; la restauration se place avant la boucle de phase et
la sauvegarde après ses contrôles de segment.

\subsection{Algorithme 3 : boucles à l'intérieur d'un segment}
\begin{algocode}{Algorithme 3 --- Metropolis, échanges et mesures dans le bon ordre}
\textbf{Entrées :} états courants, énergies, identités, parité et $S$ balayages.
\textbf{Sorties :} mêmes états mis à jour, séries et compteurs du segment.
\begin{enumerate}
\item Initialiser le générateur du noyau avec la graine reçue ; calculer
$\widetilde\beta_r=1/(Q\kb T_r)$ ; allouer $S/d_m$ lignes de mesures par
température et mettre les compteurs du segment à zéro.
\par\textit{Commentaire :} les températures sont fixes ; leurs inverses
n'ont pas besoin d'être recalculés dans la boucle des propositions.
\item \textbf{Pour} $s=1,\ldots,S$ :
\begin{enumerate}
\item \textbf{Pour} $r=0,\ldots,R-1$ : effectuer exactement $N$ appels à
l'algorithme 4 sur la configuration présente à $T_r$.
\par\textit{Commentaire :} un balayage contient $N$ propositions Au--Cu,
acceptées ou refusées. Le tirage se fait avec remise : certains sites
peuvent intervenir plusieurs fois et d'autres pas du tout.
\item \textbf{Si} $d_x>0$ et $s\bmod d_x=0$, exécuter l'algorithme 5
pour les liens de la parité courante, puis incrémenter \texttt{phase}.
\par\textit{Commentaire :} toutes les températures ont d'abord subi le
même nombre de propositions locales. $d_x=0$ désactive les échanges.
\item \textbf{Si} $s\bmod d_m=0$, calculer l'algorithme 6 pour chaque
emplacement $r$, enregistrer les neuf valeurs et \texttt{labels[r]},
puis avancer l'indice de mesure.
\par\textit{Commentaire :} si les deux échéances coïncident, la mesure
porte sur l'état \emph{après} l'échange. Elle est rangée sous $T_r$, même
si l'identité de la configuration vient de changer.
\end{enumerate}
\item Retourner les séries, les compteurs locaux, les compteurs de liens
et la valeur courante de \texttt{phase}.
\par\textit{Commentaire :} conserver cette parité au segment suivant évite
de recommencer systématiquement par les mêmes liens.
\end{enumerate}
\textbf{Source :} \texttt{kernel.py : run\_segment}.
\end{algocode}

Les mouvements locaux et les échanges préservent chacun la mesure jointe
$\Pi\propto\prod_r\exp[-E_r/(\kb T_r)]$. Leur composition dans cet ordre
préserve donc cette mesure ; il n'est pas nécessaire de prétendre que la
composition complète satisfait elle-même la balance détaillée. Les échanges
de répliques accélèrent la circulation entre températures, tandis que
l'étiquette permet d'en vérifier l'efficacité.

\subsection{Algorithme 4 : proposition Au--Cu et test de Metropolis}
\begin{algocode}{Algorithme 4 --- Une tentative locale à température fixée}
\textbf{Entrées :} occupation $s$, énergie $\mathcal E$ et $\widetilde\beta_r$.
\textbf{Sortie :} état inchangé ou état échangé avec énergie cohérente.
\begin{enumerate}
\item Tirer indépendamment $i,j$ uniformément dans $\{0,\ldots,N-1\}$.
\textbf{Tant que} $s_i=s_j$, retirer les deux indices.
\par\textit{Commentaire :} un tirage d'espèces identiques, y compris $i=j$,
n'est pas compté comme tentative. Chaque paire non orientée Au--Cu a alors
la probabilité $1/(N_{\Au}N_{\Cu})$. La composition étant conservée,
la proposition inverse a exactement la même probabilité : aucun rapport
de proposition supplémentaire n'apparaît dans Metropolis.
\item Pour le calcul local, nommer $i$ le site initialement Au et $j$ le
site initialement Cu, en permutant les noms si nécessaire. Sans modifier
l'état, calculer
\[
\begin{aligned}
a_i&=\sum_{v\in V_1(i),\,v\ne j}s_v,&
a_j&=\sum_{v\in V_1(j),\,v\ne i}s_v,\\
b_i&=\sum_{v\in V_2(i),\,v\ne j}s_v,&
b_j&=\sum_{v\in V_2(j),\,v\ne i}s_v.
\end{aligned}
\]
\textit{Commentaire :} seules les liaisons incidentes aux sites échangés
peuvent changer. Leur éventuelle liaison mutuelle reste Au--Cu, donc on
l'exclut. Un voisin commun peut contribuer aux deux sites : ce sont deux
liaisons distinctes et leurs variations doivent toutes deux être comptées.
\item Évaluer la différence entière
\[
 \Delta\mathcal E=161000(a_j-a_i)+12600(b_i-b_j).
\]
\textit{Commentaire :} déplacer l'Au de $i$ vers $j$ change le nombre de
paires Au--Au de première coquille de $a_j-a_i$. Pour les secondes coquilles,
le nombre de paires Au--Cu change de $2(b_i-b_j)$ : perdre un voisin Au
et gagner un voisin Cu ont des contributions opposées. D'où
$12600=2\times6300$. La formule est exacte à composition fixée et ne
requiert que deux listes de 18 voisins.
\item \textbf{Si} $\Delta\mathcal E\le0$, accepter. \textbf{Sinon}, tirer
$u$ uniforme sur $[0,1)$ et accepter si
\[
u<\exp(-\widetilde\beta_r\Delta\mathcal E).
\]
\textit{Commentaire :} les mouvements qui abaissent l'énergie et ceux
d'énergie égale sont toujours acceptés. Un mouvement défavorable reste
possible avec le poids de Boltzmann ; sans cela, le calcul chercherait
seulement un minimum d'énergie au lieu d'échantillonner une température.
Le test conditionnel évite de calculer l'exponentielle pour $\Delta\mathcal E\le0$.
\item \textbf{Si accepté}, permuter simultanément $s_i,s_j$, effectuer
$\mathcal E\leftarrow\mathcal E+\Delta\mathcal E$ et incrémenter le compteur.
\textbf{Sinon}, conserver exactement l'état et l'énergie précédents.
\par\textit{Commentaire :} la mise à jour doit suivre la décision. Un refus
représente un temps de séjour dans la chaîne et compte dans le balayage ;
ne mesurer que les mouvements acceptés biaiserait les poids d'équilibre.
\end{enumerate}
\textbf{Sources :} \texttt{kernel.py : delta\_fast, run\_segment}.
\end{algocode}

\subsection{Algorithme 5 : échange de répliques voisines}
\begin{algocode}{Algorithme 5 --- Une session d'échange entre températures}
\textbf{Entrées :} les $R$ états, énergies, identités et \texttt{phase}.
\textbf{Sorties :} configurations éventuellement permutées ; températures fixes.
\begin{enumerate}
\item Poser $p=\texttt{phase}\bmod2$. \textbf{Pour} les liens
$r=p,p+2,p+4,\ldots$ tels que $r+1<R$, incrémenter le nombre de tentatives
sur le lien $(r,r+1)$.
\par\textit{Commentaire :} une session teste $(0,1),(2,3),\ldots$, la
suivante $(1,2),(3,4),\ldots$. Les liens d'une session ne partagent pas de
réplique ; l'alternance permet de parcourir toute l'échelle de températures.
\item Calculer le logarithme du rapport des poids :
\[
\ell=(\widetilde\beta_r-\widetilde\beta_{r+1})
       (\mathcal E_r-\mathcal E_{r+1}).
\]
\textit{Commentaire :} en effet,
\[
\log\frac{\exp[-\beta_r E_{r+1}-\beta_{r+1}E_r]}
{\exp[-\beta_r E_r-\beta_{r+1}E_{r+1}]}
=(\beta_r-\beta_{r+1})(E_r-E_{r+1}),
\]
avec $\beta_r=1/(\kb T_r)$. Les propositions d'échange sont symétriques.
\item Accepter si $\ell\ge0$ ; sinon tirer $u$ uniforme sur $[0,1)$ et
accepter si $\log(\max(u,10^{-300}))<\ell$.
\par\textit{Commentaire :} travailler en logarithmes évite de former une
exponentielle pouvant déborder. La borne du code évite $\log0$ ; elle ne
modifie le test usuel que pour des probabilités numériquement extrêmes.
\item Si accepté, permuter les lignes complètes de \texttt{states}, les
deux énergies et les deux identités ; incrémenter le nombre d'échanges
acceptés sur ce lien. Laisser $T_r$ et $T_{r+1}$ à leurs emplacements.
\par\textit{Commentaire :} énergie et identité doivent suivre la
configuration. Permuter seulement les occupations produirait une énergie
incorrecte pour le mouvement local suivant.
\item Après tous les liens de la session, incrémenter \texttt{phase}.
\end{enumerate}
\textbf{Source :} \texttt{kernel.py : run\_segment}, bloc d'échange.
\end{algocode}

\subsection{Algorithme 6 : observables d'une configuration instantanée}
\begin{algocode}{Algorithme 6 --- Mesure après un balayage sélectionné}
\textbf{Entrées :} $s$, énergie entière courante, voisins et sous-réseaux.\par
\textbf{Sortie :} $(E,\eta,q,\alpha_1,\alpha_2,p_0,p_1,p_2,p_3)$.
\begin{enumerate}
\item Mettre $p_\mu=0$ pour les quatre sous-réseaux et $B_1=B_2=0$.
\item \textbf{Pour chaque site} $i$ tel que $s_i=1$ :
\begin{enumerate}
\item ajouter $4/N$ à $p_{\texttt{sub[i]}}$ ;
\item ajouter $1-s_v$ à $B_1$ pour chacun des 12 premiers voisins et à
$B_2$ pour chacun des 6 seconds voisins.
\end{enumerate}
\textit{Commentaire :} chaque sous-réseau contient $N/4$ sites ; ainsi
$p_\mu=N_{\Au,\mu}/(N/4)$ est une fraction d'occupation, avec
$\sum_\mu p_\mu=1$ à cette composition. Partir seulement des sites Au
compte chaque paire mixte exactement une fois : il n'y a pas de facteur
$1/2$ à appliquer à $B_n$.
\item Calculer l'ordre signé et l'amplitude :
\[
\eta=\frac{p_0-1/4}{3/4},\qquad
q=\sqrt{\frac43\sum_{\mu=0}^3(p_\mu-1/4)^2}.
\]
\textit{Commentaire :} $\eta$ suit le sous-réseau 0 et change selon la
variante ; $q$ reconnaît les quatre variantes ordonnées sans choisir de
sous-réseau privilégié. Il faut calculer $q$ pour chaque configuration
puis moyenner : le calcul à partir des $\langle p_\mu\rangle$ pourrait
effacer un ordre réel lorsque la variante change au cours de la trajectoire.
\item Pour $n=1,2$, calculer
\[
\alpha_n=1-\frac{B_n}{(N/4)z_n(3/4)},\qquad z_1=12,\ z_2=6.
\]
\textit{Commentaire :} $B_n/[(N/4)z_n]$ est la probabilité conditionnelle
de trouver Cu autour d'Au. Diviser par $c_{\Cu}=3/4$ compare ce voisinage
à la composition globale ; on mesure ici l'ordre local.
\item Former la ligne avec $E=\mathcal E/Q$ en eV \emph{totaux}, puis
$\eta,q,\alpha_1,\alpha_2,p_0,\ldots,p_3$. Enregistrer séparément l'identité.
\par\textit{Commentaire :} l'énergie n'est pas recalculée ici : sa valeur
est tenue à jour après chaque mouvement et vérifiée après le segment.
Le parcours global des occupations sert aux observables structurales.
L'énergie par atome sera obtenue au stade des statistiques.
\end{enumerate}
\textbf{Source :} \texttt{kernel.py : sample\_observables}.
\end{algocode}

L'intervalle $d_m$ réduit le coût des parcours de tous les sites et la taille
des archives. Il ne garantit pas l'indépendance des mesures. Mesurer moins
souvent ne remplace donc pas l'estimation de l'autocorrélation. De même,
une configuration inchangée entre deux échéances doit être enregistrée à
nouveau : son temps de séjour fait partie de la mesure canonique.

\subsection{Algorithme 7 : moyennes et fluctuations à température fixe}
Notons $M=S_{\mathrm{prod}}/d_m$ le nombre de mesures par chaîne et
$\overline A_c=M^{-1}\sum_{m=1}^M A_{c,m}$ la moyenne de production à une
température donnée. Pour les variances thermiques, le code utilise
$v_c(A)=\overline{A^2}_c-\overline A_c^2$, équivalent à
\texttt{numpy.var} avec le diviseur $M$ ; les barres d'erreur sont ensuite
estimées par rééchantillonnage.

\begin{algocode}{Algorithme 7 --- Estimer les observables thermiques}
\textbf{Entrées :} séries de production $[C,M,R,9]$, températures et $N$.
\textbf{Sorties :} huit estimations par température.
\begin{enumerate}
\item \textbf{Pour chaque température} $T_r$, extraire les séries à
\emph{emplacement fixe} $r$ des $C$ chaînes.
\par\textit{Commentaire :} suivre une identité de réplique mélangerait
plusieurs températures et ne donnerait pas une moyenne canonique à $T_r$.
\item \textbf{Dans chaque chaîne}, calculer
\[
\begin{gathered}
u_c=\overline E_c/N,\quad
\overline\eta_c,\quad\overline q_c,\quad
\overline\alpha_{1,c},\quad\overline\alpha_{2,c},\\
\left(\frac{c_V}{\kb}\right)_c=\frac{v_c(E)}{N(\kb T_r)^2},\qquad
(X_q)_c=\frac{N v_c(q)}{\kb T_r},\\
(U_q)_c=1-\frac{\overline{q^4}_c}{3(\overline{q^2}_c)^2}.
\end{gathered}
\]
Si $\overline{q^2}_c=0$, enregistrer un cumulant non défini.
\par\textit{Commentaire :} $c_V$ est une variance d'énergie \emph{totale}
divisée par $N$ ; remplacer $E$ par $E/N$ dans la même formule ferait
perdre un facteur $N^2$. $X_q$ et $U_q$ exigent des moments de la série de
$q$ ; ils ne se déduisent pas de la seule moyenne $\overline q$.
\item Moyenner les huit estimations sur les chaînes :
$\widehat\theta_r=C^{-1}\sum_c\widehat\theta_{c,r}$.
\par\textit{Commentaire :} les variances sont calculées avant ce
moyennage. La variance d'un mélange de chaînes ajouterait la variance de
leurs moyennes à la fluctuation interne et pourrait transformer un défaut
d'équilibration en pic artificiel de capacité calorifique.
\item Transmettre estimations et séries à l'algorithme 8 pour les erreurs
et les diagnostics. Après analyse de toutes les températures, choisir
$T^*=T_{\arg\max_r(c_V/\kb)_r}$ et sauvegarder ses voisins sur la grille.
\par\textit{Commentaire :} cette recherche intervient après les mesures ;
aucune température de transition n'est imposée à la dynamique. Les deux
voisins décrivent une résolution de grille, pas une erreur-type sur $T^*$.
\end{enumerate}
\textbf{Sources :} \texttt{statistics.py : moments, summarize} ;
\texttt{simulation.py : campaign}.
\end{algocode}

\subsection{Algorithme 8 : autocorrélation, incertitudes et diagnostics}
\begin{algocode}{Algorithme 8 --- Analyse statistique des séries corrélées}
\textbf{Entrées :} séries de production à $T_r$, $B$ tirages bootstrap.
\textbf{Sorties :} erreurs-types, temps de corrélation et critères de suivi.
\begin{enumerate}
\item Pour chaque chaîne et pour $A=E,q$, centrer la série. Calculer par
transformée de Fourier sa covariance aux retards $\ell$, en divisant par
$M-\ell$, puis normaliser par la covariance à retard nul pour obtenir
$\rho_A(\ell)$.
\par\textit{Commentaire :} le remplissage de zéros de la transformée évite
une corrélation circulaire parasite ; la FFT accélère le calcul de tous les
retards. Une série de moins de huit valeurs ou de variance numérique
quasi nulle reçoit $\tau=1/2$ dans le code ; ce retour conventionnel ne
constitue pas une preuve de mélange si la trajectoire est figée.
\item Accumuler les paires
$P_j=\rho_A(2j-1)+\rho_A(2j)$ tant qu'elles sont positives. Arrêter à la
première paire non positive ou selon la fenêtre du code,
$2j-1>10[1/2+\sum_{h\le j}P_h]$. Poser
\[
\tau_{c,A}=\max\left(\frac12,\frac12+\sum_{j\ \text{retenus}}P_j\right),
\qquad \tau_{\max}=\max_{c,A\in\{E,q\}}\tau_{c,A}.
\]
\textit{Commentaire :} les grands retards sont bruyants ; sommer sans
troncature rendrait l'estimation instable. $\tau$ est exprimé en
\emph{mesures enregistrées} ; multiplier par $d_m$ pour l'exprimer en balayages.
\item Choisir exactement comme dans \texttt{summarize} la longueur de bloc
\[
b=\min\!\left(\max(2,\lceil2\tau_{\max}\rceil),
                \max(2,\lfloor M/4\rfloor)\right).
\]
\textit{Commentaire :} on vise des blocs au moins aussi longs que
$2\tau_{\max}$ pour conserver la dépendance temporelle. Le plafond lié à
$M/4$ limite cette longueur pour les séries courtes ; si ce plafond devient
actif, la longueur souhaitée n'est plus garantie.
\item \textbf{Pour} chaque tirage bootstrap $h=1,\ldots,B$ et chaque chaîne :
tirer $\lceil M/b\rceil$ débuts de blocs uniformément, construire les
indices consécutifs modulo $M$, concaténer et garder les $M$ premiers.
Rééchantillonner ensemble les neuf colonnes avec ces indices ; recalculer
les huit fonctions de l'algorithme 7 puis moyenner sur les chaînes.
\par\textit{Commentaire :} les blocs circulaires évitent de sous-représenter
les bords. Recalculer les variances et le cumulant à chaque tirage transmet
l'incertitude à ces fonctions non linéaires. Le rééchantillonnage est fait
séparément dans chaque chaîne et à chaque température.
\item Pour chaque estimation $\theta$, retenir
\[
\mathrm{SE}(\theta)=\max\!\left[
\mathrm{sd}_{h}(\widehat\theta^{*(h)}),\,
\frac{\mathrm{sd}_{c}(\widehat\theta_c)}{\sqrt C}\right],
\]
avec des écarts-types à diviseur effectif nombre de valeurs moins un.
\par\textit{Commentaire :} le second terme empêche d'afficher une erreur
plus petite que la dispersion observée entre chaînes. Ce maximum reste
une estimation d'erreur-type, sans corriger une éventuelle métastabilité
commune aux deux chaînes.
\item Calculer l'effectif utile rapporté
\[
\mathrm{ESS}_{\min}=\min_{A\in\{E,q\}}
       \sum_{c=1}^C\frac{M}{2\tau_{c,A}},
\]
ainsi que $\widehat R_E$, $\widehat R_q$ et le nombre de blocs
$\lfloor M/b\rfloor$ par chaîne.
\item Comparer les deux demi-séries d'énergie de chaque chaîne et retenir
\[
z_{\max}=\max_c
\frac{|\overline E_{c,\text{première moitié}}
       -\overline E_{c,\text{seconde moitié}}|}
{\max[s_c(E)\sqrt{8\tau_{\max}/M},10^{-12}]},
\]
où $s_c(E)$ est l'écart-type à diviseur $M-1$. Déclarer le critère local
satisfait si simultanément
\[
\mathrm{ESS}_{\min}\ge200,\quad\widehat R_E<1.05,\quad
\widehat R_q<1.05,\quad\lfloor M/b\rfloor\ge20,\quad z_{\max}<3.
\]
\textit{Commentaire :} ces seuils signalent les séries à renforcer ; ils
ne garantissent pas que tous les bassins métastables ont été visités.
\end{enumerate}
\textbf{Sources :} \texttt{statistics.py : tau\_int, split\_rhat, summarize}.
\end{algocode}

Pour préciser le calcul de $\widehat R$, le code partage chaque chaîne en
deux moitiés de longueur $n_h=\lfloor M/2\rfloor$, puis normalise les rangs
de leurs valeurs par
\[
z=\Phi^{-1}\!\left(\frac{\mathrm{rang}-3/8}{2Cn_h+1/4}\right).
\]
Il calcule $W$, la moyenne des variances internes, et
$B_{\mathrm{inter}}=n_h\,\mathrm{Var}(\overline z_{\text{demi-chaîne}})$,
puis
\[
\widehat R=\max\!\left(1,
\sqrt{\frac{(n_h-1)W/n_h+B_{\mathrm{inter}}/n_h}{W}}\right).
\]
Le résultat retenu est le maximum de cette version et de celle obtenue
en remplaçant les valeurs initiales par leurs écarts absolus à la médiane
globale, avant de recalculer les rangs. Cette seconde version recherche des
différences de dispersion. Les cas dégénérés $W\simeq0$ et les séries trop
courtes sont traités séparément par \texttt{split\_rhat}. L'ordre signé
$\eta$ n'est pas utilisé ici car deux variantes ordonnées peuvent avoir
des valeurs signées différentes tout en décrivant le même degré d'ordre.

Les identités fournissent un autre contrôle, calculé sur la production :
pour chaque étiquette, un automate attend successivement une présence au
froid, au chaud, puis au froid, et compte alors un aller-retour. Ce comptage
utilise les seules dates enregistrées et peut manquer des visites entre
deux mesures. Les taux locaux sont
$N_{\text{acceptés},r}/(NS_{\mathrm{prod}})$ et les taux d'échange sont
$N_{\text{échanges acceptés},r}/N_{\mathrm{tentatives},r}$ par lien.
Ces contrôles de circulation complètent les critères de séries temporelles.

\subsection{Chronologie concrète avec les paramètres de production}
Le fichier \texttt{configs/production.json}, également recopié dans
\texttt{summary.json}, fixe $N=864$, $R=46$, $C=2$,
$S_{\mathrm{eq}}=6000$, $S_{\mathrm{prod}}=30000$, $K=500$,
$d_x=5$, $d_m=10$ et $B=250$. Un segment contient donc 500 balayages,
100 sessions d'échange et 50 dates de mesure. Une session ne teste
qu'une parité de liens : un lien donné est sollicité toutes les deux
sessions, soit tous les 10 balayages dans cette campagne.

\begin{center}\small
\begin{tabular}{r p{.66\linewidth}}
\toprule
Balayage local & Opérations après les $N$ propositions à chacune des $R$ températures\\
\midrule
1 à 4 & Pas d'échange et pas de mesure enregistrée.\\
5 & Tester $(0,1),(2,3),\ldots,(44,45)$ ; aucune mesure.\\
6 à 9 & Poursuivre les mouvements locaux.\\
10 & Tester $(1,2),(3,4),\ldots,(43,44)$ ; mesurer ensuite les 46 états.\\
15 & Reprendre la première parité de liens.\\
20 & Tester la seconde parité puis mesurer.\\
500 & Échanges, dernière mesure du segment, puis vérification globale des énergies et de la composition.\\
\bottomrule
\end{tabular}
\end{center}

Ce tableau décrit le premier segment ; les suivants poursuivent la valeur
de \texttt{phase}. Ici les 100 sessions par segment ramènent à la même
parité, mais le code conserve explicitement cette variable pour les autres
configurations admissibles. La production comporte 60 segments par chaîne
et $M=3000$ mesures par température et par chaîne. L'équilibration comporte
12 segments et 600 mesures diagnostiques par température et par chaîne,
écartées des estimations thermiques.

Au total, hors tirages rejetés d'espèces identiques, la campagne comporte
\[
C R N(S_{\mathrm{eq}}+S_{\mathrm{prod}})
=2\times46\times864\times36000
=2\,861\,568\,000
\]
propositions Au--Cu. Ce total compte les tests de Metropolis, et non le
nombre de mouvements finalement acceptés.

\subsection{Validation finale, champ moyen et coût du calcul}
Après la campagne, \texttt{validation.py} vérifie notamment, pour toutes
les mesures archivées, les identités sur $p_\mu$, $\eta$ et $q$, puis
reconstruit l'énergie à partir des paramètres de Cowley :
\[
B_1=(1-\alpha_1)(N/4)12(3/4),\quad
A_1=\frac{12(N/4)-B_1}{2},\quad
B_2=(1-\alpha_2)(N/4)6(3/4),
\]
\[
E_{\mathrm{reconstruit}}=-3.813N+0.161A_1+0.0063B_2.
\]
Cette comparaison indépendante associe les observables structurales à
l'énergie effectivement mesurée. La validation compare aussi les moyennes
d'énergie du CSV aux archives et contrôle les bornes analytiques ainsi que
les régimes froid et chaud. Le lanceur exécute auparavant les tests
physiques et statistiques ; un échec interrompt le calcul.

La référence de champ moyen est obtenue séparément dans
\texttt{meanfield.py}. L'algorithme de Brent recherche une racine de
$h(\eta)=\eta g'(\eta)-2[g(\eta)-g(0)]$ sur
$[0.01,0.999999]$, puis calcule
$T_{\mathrm{MF}}=2A\eta/[\kb g'(\eta)]$ et
$T_{\mathrm{sp}}=2A/\kb$. L'intervalle évite la solution triviale
$\eta=0$ et la singularité de $g'$ à $\eta=1$ ; les termes entropiques
$p\log p$ utilisent une fonction qui applique $0\log0=0$.
Ce calcul fournit une comparaison analytique sans modifier les probabilités
de transition Monte-Carlo.

À nombre de voisins fixé, une différence d'énergie locale coûte $O(1)$
et un balayage de toutes les températures coûte $O(RN)$ en moyenne.
Le contrôle global d'énergie coûte également $O(RN)$, une seule fois par
segment ; le calcul des observables coûte $O(RN)$ à chaque date de mesure.
La précalculation des voisins et la compilation Numba réduisent le travail
dans les boucles de propositions. Les configurations et voisinages
occupent $O(RN+18N)$, tandis que les neuf séries de production occupent
$O(9CMR)$ valeurs, auxquelles s'ajoutent identités et équilibration.
Le bootstrap, effectué après la dynamique, réévalue les statistiques sur
des blocs pour propager les corrélations temporelles dans les erreurs.

\clearpage
\subsection{Version en pseudocode}
\label{sec:pseudocode}
Les blocs suivants donnent une lecture directement transposable en code
des algorithmes précédents. L'indentation et les mots-clés
\texttt{POUR}, \texttt{TANT QUE}, \texttt{SI} indiquent les boucles et les
conditions ; \texttt{<-} désigne une affectation et \texttt{\#} un commentaire.
Les indices commencent à zéro. Les tableaux passés à une procédure sont
modifiés sur place, sauf les données explicitement retournées.

On conserve les conventions de l'annexe : \texttt{Q = 1000000},
\texttt{Eint} est une énergie totale en micro-eV entiers, \texttt{E} une
énergie totale en eV, \texttt{dm} l'intervalle de mesure, \texttt{dx}
l'intervalle d'échange et \texttt{K} la taille maximale d'un segment.
\texttt{UNIFORME()} tire un réel dans $[0,1)$ et
\texttt{ENTIER(a,b)} un entier uniforme dans $\{a,\ldots,b\}$.
Chaque commentaire signale le rôle de l'étape dans le calcul.

\begin{algocode}{Pseudocode P1 --- Géométrie et énergie de référence}
\small
\begin{verbatim}
FONCTION GEOMETRIE(L):
    EXIGER L entier ET L >= 3
    base <- [(0,0,0), (0,1,1), (1,0,1), (1,1,0)]
    xyz, sub <- listes vides
    POUR chaque maille (a,b,c) dans {0,...,L-1}^3:
        POUR mu = 0,...,3:
            AJOUTER 2*(a,b,c) + base[mu] A xyz
            AJOUTER mu A sub
        FIN POUR
    FIN POUR
    N <- longueur(xyz)
    index <- dictionnaire {coordonnees : indice du site}
    D1 <- deplacements entiers de norme carree 2
    D2 <- [(+/-2,0,0), (0,+/-2,0), (0,0,+/-2)]
    POUR i = 0,...,N-1:
        V1[i] <- [index[(xyz[i]+d) modulo (2*L)] pour d dans D1]
        V2[i] <- [index[(xyz[i]+d) modulo (2*L)] pour d dans D2]
    FIN POUR
    RETOURNER xyz, sub, V1, V2  # Voisinages calcules une fois

FONCTION ENERGIE_REFERENCE(s):
    total <- 0
    POUR i = 0,...,N-1:
        POUR v dans V1[i]: total <- total + Jint[s[i],s[v]]
        POUR v dans V2[i]: total <- total + Iint[s[i],s[v]]
    FIN POUR
    RETOURNER total / 2        # Chaque paire est vue deux fois
\end{verbatim}
\end{algocode}

\texttt{Jint} et \texttt{Iint} sont les matrices de paires multipliées par
\texttt{Q}. La division finale donne un entier exact. Les signes
\texttt{+/-} du pseudocode signifient que les deux signes sont énumérés.
La taille $N=4L^3$ est déduite du nombre de sites construits.

\begin{algocode}{Pseudocode P2 --- Campagne : chaînes, phases et segments}
\small
\begin{verbatim}
VERIFIER_PARAMETRES(cfg)             # Multiples de dm et dx
xyz, sub, V1, V2 <- GEOMETRIE(L)
N <- longueur(xyz)
archives <- liste vide

POUR c = 0,...,C-1:
    rng <- GENERATEUR(graine[c])     # Generateur de cette chaine
    POUR r = 0,...,R-1:
        SI c = 0 OU T[r] < Tdepart:
            s[r] <- (sub = (2*c) modulo 4)
        SINON:
            s[r] <- N/4 sites Au tires sans remise avec rng
        FIN SI
        Eint[r] <- ENERGIE_REFERENCE(s[r])
    FIN POUR
    labels <- [0,...,R-1]
    phase <- 0                      # Parite des liens a echanger
    burn, prod, ids <- listes vides
    acc_prod, ok_prod, essais_prod <- tableaux nuls

    POUR (mode, longueur) dans [(equilibration, Seq),
                               (production, Sprod)]:
        fait <- 0
        TANT QUE fait < longueur:
            S <- min(K, longueur - fait)
            g <- TIRER_GRAINE(rng)  # Nouvelle graine de segment
            rec, lab, acc, ok, essais, phase <- SEGMENT(
                s, Eint, labels, phase, S, g)
            POUR r = 0,...,R-1:
                EXIGER ENERGIE_REFERENCE(s[r]) = Eint[r]
                EXIGER somme(s[r]) = N/4
                EXIGER Eint[r] >= ARRONDIR(-3.813*N*Q)
            FIN POUR                # Controle apres le segment
            SI mode = equilibration:
                CONCATENER rec A burn
            SINON:
                CONCATENER rec A prod ET lab A ids
                acc_prod <- acc_prod + acc
                ok_prod <- ok_prod + ok
                essais_prod <- essais_prod + essais
            FIN SI
            fait <- fait + S
        FIN TANT QUE
    FIN POUR                        # Pas de remise a zero des etats
    SAUVEGARDER(prod, burn, ids, s, T, xyz, sub)
    AJOUTER prod A archives
    CALCULER_ACCEPTATIONS_ET_ALLERS_RETOURS(ids,
        acc_prod, ok_prod, essais_prod)
FIN POUR

tableau <- ANALYSER(archives, T, N)
Tetoile <- temperature du maximum de tableau.cv
SAUVEGARDER(tableau, Tetoile, voisins_de_grille(Tetoile))
VALIDER_ARCHIVES_ET_CSV()
PRODUIRE_FIGURES_ET_RAPPORT()
\end{verbatim}
\end{algocode}

\texttt{burn} conserve seulement les historiques d'équilibration.
\texttt{ANALYSER} reçoit uniquement \texttt{prod}, donc les mesures
initiales sont exclues des estimations thermiques. Les fonctions de
sauvegarde et de contrôle des allers-retours correspondent à celles
décrites dans l'algorithme 2 et après l'algorithme 8.

\begin{algocode}{Pseudocode P3 --- Un segment de balayages et d'échanges}
\small
\begin{verbatim}
PROCEDURE SEGMENT(s, Eint, labels, phase, S, graine):
    INITIALISER_ALEATOIRE_NOYAU(graine)
    beta[r] <- 1 / (Q*kB*T[r]) pour chaque r
    rec, lab <- tableaux de S/dm dates
    acc, ok, essais <- tableaux nuls
    m <- 0

    POUR balayage = 1,...,S:
        POUR r = 0,...,R-1:
            POUR tentative = 1,...,N:
                accepte <- METROPOLIS(s[r], Eint[r], beta[r])
                SI accepte: acc[r] <- acc[r] + 1
            FIN POUR                # N propositions, refus inclus
        FIN POUR

        SI dx > 0 ET balayage modulo dx = 0:
            POUR r = phase modulo 2,...,R-2 PAR PAS DE 2:
                essais[r] <- essais[r] + 1
                ell <- (beta[r]-beta[r+1])*(Eint[r]-Eint[r+1])
                SI ell >= 0:
                    accepter <- VRAI
                SINON:
                    u <- UNIFORME()
                    accepter <- log(max(u, 1e-300)) < ell
                FIN SI
                SI accepter:
                    PERMUTER s[r], s[r+1]       # Lignes completes
                    PERMUTER Eint[r], Eint[r+1]
                    PERMUTER labels[r], labels[r+1]
                    ok[r] <- ok[r] + 1
                FIN SI
            FIN POUR
            phase <- phase + 1      # Alterner les liens testes
        FIN SI

        SI balayage modulo dm = 0:
            POUR r = 0,...,R-1:
                rec[m,r] <- OBSERVABLES(s[r], Eint[r])
                lab[m,r] <- labels[r]
            FIN POUR                # Mesure APRES les echanges
            m <- m + 1
        FIN SI
    FIN POUR
    RETOURNER rec, lab, acc, ok, essais, phase
\end{verbatim}
\end{algocode}

Les températures ne sont jamais permutées : \texttt{rec[m,r]} désigne
toujours $T_r$. Le compteur de balayages repart de 1 à chaque appel, mais
les multiples exigés dans P2 assurent la continuité du calendrier des
mesures et des échanges. La variable \texttt{phase} est retournée pour
conserver l'alternance entre appels.

\begin{algocode}{Pseudocode P4 --- Variation locale et décision de Metropolis}
\small
\begin{verbatim}
FONCTION DELTA_LOCAL(s, i, j):
    SI s[i] = 0: PERMUTER les noms i, j
    # Ici i porte Au et j porte Cu ; s est encore inchange.
    ai <- somme(s[v] pour v dans V1[i] si v != j)
    aj <- somme(s[v] pour v dans V1[j] si v != i)
    bi <- somme(s[v] pour v dans V2[i] si v != j)
    bj <- somme(s[v] pour v dans V2[j] si v != i)
    RETOURNER 161000*(aj-ai) + 12600*(bi-bj)

PROCEDURE METROPOLIS(s, Eint, beta):
    REPETER:
        i <- ENTIER(0,N-1)
        j <- ENTIER(0,N-1)
    JUSQU'A s[i] != s[j]             # Proposer uniquement Au-Cu
    dE <- DELTA_LOCAL(s, i, j)       # Energie entiere en micro-eV
    SI dE <= 0:
        accepter <- VRAI
    SINON:
        accepter <- UNIFORME() < exp(-beta*dE)
    FIN SI
    SI accepter:
        PERMUTER s[i], s[j]          # Composition conservee
        Eint <- Eint + dE           # Mise a jour apres decision
    FIN SI                          # Un refus garde l'ancien etat
    RETOURNER accepter
\end{verbatim}
\end{algocode}

\texttt{Eint} est ici un argument modifiable : sa nouvelle valeur est
transmise à l'appelant. Dans \texttt{DELTA\_LOCAL}, permuter les noms
\texttt{i,j} ne modifie pas les occupations et ne change pas les indices
de l'appelant. La liaison mutuelle des sites, si elle existe, est exclue
des quatre sommes car elle reste Au--Cu après l'échange.

\begin{algocode}{Pseudocode P5 --- Calcul des neuf observables instantanées}
\small
\begin{verbatim}
FONCTION OBSERVABLES(s, Eint):
    p <- [0,0,0,0]
    B1, B2 <- 0, 0
    POUR i = 0,...,N-1:
        SI s[i] = 1:
            p[sub[i]] <- p[sub[i]] + 4/N
            POUR v dans V1[i]: B1 <- B1 + (1-s[v])
            POUR v dans V2[i]: B2 <- B2 + (1-s[v])
        FIN SI
    FIN POUR                        # Chaque paire mixte vue une fois
    eta <- (p[0]-1/4) / (3/4)
    q <- sqrt((4/3)*somme((p[mu]-1/4)^2 pour mu=0,...,3))
    alpha1 <- 1 - B1 / ((N/4)*12*(3/4))
    alpha2 <- 1 - B2 / ((N/4)*6*(3/4))
    E <- Eint / Q                   # Energie TOTALE en eV
    RETOURNER [E, eta, q, alpha1, alpha2, p[0],p[1],p[2],p[3]]
\end{verbatim}
\end{algocode}

\begin{algocode}{Pseudocode P6 --- Moyennes thermiques et bootstrap par blocs}
\small
\begin{verbatim}
FONCTION MOMENTS(serie, T, N):
    E, eta, q, alpha1, alpha2 <- cinq premieres colonnes
    u  <- moyenne(E)/N
    cv <- variance(E) / (N*(kB*T)^2) # Ici cv signifie cV/kB
    Xq <- N*variance(q) / (kB*T)
    SI moyenne(q^2) > 0:
        Uq <- 1 - moyenne(q^4)/(3*moyenne(q^2)^2)
    SINON:
        Uq <- NON_DEFINI
    FIN SI
    RETOURNER [u, moyenne(eta), moyenne(q), moyenne(alpha1),
               moyenne(alpha2), cv, Xq, Uq]

FONCTION ANALYSER(archives, T, N):
    rng_boot <- GENERATEUR(318)
    POUR r = 0,...,R-1:
        a <- archives[:,:,r,:]      # Temperature fixe, production
        M <- nombre de mesures par chaine
        POUR c = 0,...,C-1:
            tau[c,E] <- TAU_INT(a[c,:,E])
            tau[c,q] <- TAU_INT(a[c,:,q])
            theta[c] <- MOMENTS(a[c], T[r], N)
        FIN POUR
        tau_max <- maximum(tau)
        b <- min(max(2, plafond(2*tau_max)), max(2, plancher(M/4)))
        estimation <- moyenne(theta sur les chaines)

        POUR h = 1,...,B:
            POUR c = 0,...,C-1:
                debuts <- plafond(M/b) entiers dans [0,M-1]
                          tires avec rng_boot
                ix <- CONCATENER((debut + [0,...,b-1]) modulo M
                                  pour chaque debut dans debuts)
                ix <- les M premiers indices
                tirage[c] <- MOMENTS(a[c,ix,:], T[r], N)
            FIN POUR                # Blocs tires par chaine
            boot[h] <- moyenne(tirage sur les chaines)
        FIN POUR
        erreur <- max_composante(ecart_type(boot),
                                 ecart_type(theta)/sqrt(C))
        ess <- min_sur_E_et_q(somme_sur_c(M/(2*tau[c,A])))
        rE <- SPLIT_RHAT(a[:,:,E])
        rq <- SPLIT_RHAT(a[:,:,q])
        z <- ECART_DEMI_SERIES_ENERGIE(a, tau_max)
        diagnostic <- (ess >= 200 ET rE < 1.05 ET rq < 1.05
                       ET plancher(M/b) >= 20 ET z < 3)
        AJOUTER T[r], estimation, erreur, tau_max, ess,
                rE, rq, b, z, diagnostic AU tableau
    FIN POUR
    RETOURNER tableau
\end{verbatim}
\end{algocode}

Dans P6, \texttt{variance} utilise le diviseur $M$, tandis que les
\texttt{ecart\_type} du bootstrap et de la dispersion entre chaînes
utilisent le diviseur nombre de valeurs moins un. Le maximum des erreurs
est pris séparément pour chacune des huit observables. Les routines
\texttt{TAU\_INT}, \texttt{SPLIT\_RHAT} et
\texttt{ECART\_DEMI\_SERIES\_ENERGIE} suivent les formules et les
traitements des cas dégénérés explicités dans l'algorithme 8. Les énergies
et $q$ sont des colonnes de la série ; le symbole \texttt{A} de la ligne
\texttt{ess} parcourt ces deux colonnes.

\endgroup
\begin{thebibliography}{9}
\bibitem{metropolis} N. Metropolis, A. W. Rosenbluth, M. N. Rosenbluth,
A. H. Teller et E. Teller, \emph{Equation of State Calculations by Fast
Computing Machines}, J. Chem. Phys. \textbf{21}, 1087--1092 (1953).
\href{https://doi.org/10.1063/1.1699114}{doi:10.1063/1.1699114}.
\bibitem{hukushima} K. Hukushima et K. Nemoto, \emph{Exchange Monte Carlo
Method and Application to Spin Glass Simulations}, J. Phys. Soc. Jpn.
\textbf{65}, 1604--1608 (1996).
\href{https://doi.org/10.1143/JPSJ.65.1604}{doi:10.1143/JPSJ.65.1604}.
\bibitem{binder} K. Binder, \emph{Theory of first-order phase transitions},
Rep. Prog. Phys. \textbf{50}, 783 (1987).
\href{https://doi.org/10.1088/0034-4885/50/7/001}{doi:10.1088/0034-4885/50/7/001}.
\bibitem{benisek} A. Benisek et E. Dachs, \emph{The vibrational and configurational
entropy of disordering in Cu$_3$Au}, Journal of Alloys and Compounds
\textbf{632}, 585--590 (2015).
\href{https://doi.org/10.1016/j.jallcom.2014.12.215}{doi:10.1016/j.jallcom.2014.12.215}.
\bibitem{owen} L. R. Owen, H. Y. Playford, H. J. Stone et M. G. Tucker,
\emph{Analysis of short-range order in Cu$_3$Au using X-ray pair distribution functions}.
\href{https://doi.org/10.17863/CAM.38392}{Manuscrit accepté, dépôt de Cambridge}.
\bibitem{kessler} M. Kessler, W. Dieterich et A. Majhofer,
\emph{Ordering kinetics in an fcc A$_3$B binary alloy model: Monte Carlo studies},
Phys. Rev. B \textbf{67}, 134201 (2003).
\href{https://arxiv.org/abs/cond-mat/0207611}{Prépublication et protocole de sauts atome--lacune}.
\bibitem{paras} J. Paras et A. Allanore, \emph{Contribution of electronic
entropy to the order-disorder transition of Cu$_3$Au}, Phys. Rev. Research
\textbf{3}, 023239 (2021).
\href{https://doi.org/10.1103/PhysRevResearch.3.023239}{doi:10.1103/PhysRevResearch.3.023239}.
\end{thebibliography}
\end{document}
